\documentclass[11pt]{article}
\usepackage[T1]{fontenc}
\usepackage{lmodern}
\usepackage[margin=1in]{geometry}
\usepackage{amsmath,amssymb,amsthm,booktabs,microtype}
\usepackage[hidelinks]{hyperref}
\newtheorem{theorem}{Theorem}[section]
\newtheorem{lemma}[theorem]{Lemma}
\newtheorem{proposition}[theorem]{Proposition}
\newtheorem{corollary}[theorem]{Corollary}
\theoremstyle{remark}
\newtheorem{remark}[theorem]{Remark}
\title{Nonintegral Laplacian spectra of ideal intersection graphs}
\author{Working manuscript}
\date{October 4, 2026}
\begin{document}
\maketitle
\begin{abstract}
We study Laplacian integrality for the graph of nonzero proper ideals of
$\mathbb Z_n$, where distinct ideals are adjacent when their intersection
is nonzero. For squarefree composite $n$, the graph is Laplacian integral
exactly when $n$ has two prime factors. For three distinct primes, we prove
nonintegrality for every repeated-exponent triple $(a,a,b)$ with $a,b\geq1$
and every exponent triple with minimum entry at most seven. For ordered
$2\leq a\leq b\leq c$, we also prove nonintegrality whenever
$c\geq4a^2-2a$.
A symmetric Schur complement supplies a positive quotient root in $(0,3)$
for minimum exponent at least four and proves further nonintegral families,
including every exponent span at most three and an unbounded balanced region.
Arithmetic arguments settle every common exponent divisor at least three,
every all-odd triple, and every common $2$-adic valuation.
Shifted root distributions settle every permutation of $(3,3,2)$ and
$(3,0,0)$ modulo four. A combined mixed-parity theorem includes these exclusions and, for residue
roles $(1,3,2)$ modulo four, combines the sum congruence
$a+c\equiv b(b+2)\pmod{128}$, a derivative restriction on $b$, and a
binary-cubic parity condition. For minimum at least four, it also requires
the endpoint-two equation and a linear bound on the maximum exponent.
The same linear bound applies to all-even
triples and mixed pattern $(3,3,0)$ modulo four.
The repeated-exponent proof combines discriminant arithmetic with a
consecutive-integer interval for a symmetric cubic. Exact finite certificates
exhaust the exceptional divisor levels, the minimum-three region and the
cutoff regions at minima four through seven; an endpoint certificate also
settles every triple with second-smallest exponent at most fifteen.
The remaining unequal-exponent and nonsquarefree higher-prime cases stay open.
\end{abstract}
\section{Introduction}

The ideal intersection graph associates a graph to the intersection
structure of the nonzero proper ideals of a ring~\cite{intersection}.
For $\mathbb Z_n$, its isomorphism type depends on the exponents in the
prime factorization of $n$. We ask which of these graphs have integer
Laplacian eigenvalues.

We prove a complete classification for squarefree composite integers
(Theorem~\ref{thm:squarefree}), an explicit quadratic obstruction for
$(1,1,k)$ (Theorem~\ref{thm:pqrk}), and a linear cutoff for repeated
exponents (Theorem~\ref{thm:second}). The cutoff combines an arithmetic
restriction on the antisymmetric block with the symmetric cubic obstruction
for $b=(a-1)(2a-1)$ proved in Theorem~\ref{thm:boundary}.
An exact divisor criterion isolates every possible square discriminant for
fixed $a$, leading to nonintegrality for all $b$ when $a=2,3,4,5,6$
(Theorem~\ref{thm:small} and Corollary~\ref{cor:five-six}).

The graph and lifting arguments are given explicitly, so the proofs do not
require a complete spectral decomposition. Section~\ref{sec:verification}
describes the exact computations supporting the displayed formulas and
finite certificates. The classification for arbitrary exponent vectors
remains open.

\section{The graph and its invariant subspaces}
\label{sec:graph}

Write $n=\prod_{i=1}^t p_i^{k_i}$, where the primes are distinct and
$k_i\geq1$. Denote by $G_n$ the graph whose vertices are the nonzero
proper ideals of $\mathbb Z_n$, with distinct ideals adjacent if their
intersection is nonzero. We use the combinatorial Laplacian $L=D-A$;
the graph is \emph{Laplacian integral} when every eigenvalue of $L$ is an
integer, including when the graph is disconnected.

The ideal $(\prod_i p_i^{\alpha_i})$ is encoded by
$0\leq\alpha_i\leq k_i$. We exclude $(0,\ldots,0)$, which encodes the
whole ring, and $(k_1,\ldots,k_t)$, which encodes the zero ideal.
Ideal intersection is generated by the least common multiple. Therefore
two distinct vertices $\alpha,\beta$ are adjacent exactly when
\[
\max(\alpha_i,\beta_i)<k_i\quad\hbox{for some }i.
\]
Equivalently, putting $S(\alpha)=\{i:\alpha_i<k_i\}$, adjacency is
$S(\alpha)\cap S(\beta)\ne\varnothing$.

For nonempty $S\subseteq[t]$, the support class $\mathcal C_S$ is a
clique of size $\prod_{i\in S}k_i$, except that the full support has
size $u=\prod_i k_i-1$. Edges between distinct support classes are complete
or absent according as their supports intersect or are disjoint. The full
support consists of universal vertices. Remove them and call the resulting
graph $H$; thus $G_n=K_u\vee H$, with the evident interpretation for $u=0$.

\begin{lemma}[Lifting across the universal vertices]
\label{lem:join}
If $L(H)v=\lambda v$ and $\lambda\ne0$, extension of $v$ by zero on
the $u$ universal vertices is an eigenvector of $L(G_n)$ with eigenvalue
$\lambda+u$.
\end{lemma}
\begin{proof}
Since $L(H)$ is symmetric with zero row sums, $\lambda\ne0$ implies
$\sum v=0$. At a vertex of $H$, its degree increases by $u$ and all
new neighbor values are zero. At a universal vertex the Laplacian action
is $-\sum v=0$. These are exactly the claimed eigenvector equations.
\end{proof}

For $(k_1,k_2,k_3)=(a,a,b)$, $H$ has six support classes of sizes
$a,a,b,a^2,ab,ab$ on $1,2,3,12,13,23$, respectively. The partition into
\[
X=\mathcal C_1\cup\mathcal C_2,\quad Y=\mathcal C_3,\quad
Z=\mathcal C_{12},\quad W=\mathcal C_{13}\cup\mathcal C_{23}
\]
is equitable, and on functions constant on these cells the Laplacian acts by
\begin{equation}
\label{eq:symmetric}
B=\begin{pmatrix}
a^2+ab&0&-a^2&-ab\\
0&2ab&0&-2ab\\
-2a&0&2a+2ab&-2ab\\
-a&-b&-a^2&a^2+a+b
\end{pmatrix}.
\end{equation}
For example, an $X$ vertex has $a^2$ neighbors in $Z$ and $ab$ in $W$,
while a $W$ vertex has $a,b,a^2$ neighbors in $X,Y,Z$.
Internal neighbors cancel for cell-constant functions. The cell sizes are
$2a,b,a^2,2ab$, and multiplying $B$ on the left by their diagonal matrix
gives a symmetric matrix. Thus $B$ is similar to a real symmetric matrix.

The functions taking values $s,-s$ on $\mathcal C_1,\mathcal C_2$,
$t,-t$ on $\mathcal C_{13},\mathcal C_{23}$, and zero elsewhere
form a second invariant subspace, with matrix
\begin{equation}
\label{eq:antisymmetric}
B_-=\begin{pmatrix}a^2+ab&-ab\\-a&a^2+a+b+2ab\end{pmatrix}.
\end{equation}
These functions have sum zero. Their eigenvalues therefore also lift by
$u=a^2b-1$, directly by the proof of Lemma~\ref{lem:join}.

Expanding~\eqref{eq:symmetric} gives $\det(xI-B)=x f_{a,b}(x)$, where
\begin{align}
f_{a,b}(x)&=x^3-E_1x^2+E_2x-E_3,\label{eq:cubic}\\
E_1&=2a^2+5ab+3a+b,\notag\\
E_2&=a^4+7a^3b+3a^3+8a^2b^2+11a^2b+2a^2+3ab^2+2ab,\notag\\
E_3&=2a^2b(a+b+1)(a^2+2ab+2a+b).\notag
\end{align}
Its roots are real and nonzero: reality follows from the similarity above,
and $E_3>0$. Every root $\lambda$ gives an eigenvalue $\lambda+u$ of
$G_n$. A noninteger root of $f_{a,b}$ therefore proves nonintegrality.

\section{Squarefree integers}

For squarefree integers the support classes have size one, so the graph
is described directly by subsets of the prime factors. Alternating tensors
provide invariant subspaces for this subset action~\cite{fultonharris} and
give the complete squarefree classification.

\begin{theorem}
\label{thm:squarefree}
For squarefree composite $n$, the graph $G_n$ is Laplacian integral if
and only if $n$ has exactly two prime factors.
\end{theorem}
\begin{proof}
Let $t$ be the number of prime factors. For $t=2$ the graph consists of
the two ideals $(p)$ and $(q)$, with no edge, so $L=0$. Assume $t\geq3$.
The vertices are the nonempty proper subsets of $[t]$, adjacent when they
intersect. On the tensor space encoding all subsets put
\[
E=\begin{pmatrix}2&0\\0&1\end{pmatrix},\quad
F=\begin{pmatrix}1&1\\1&0\end{pmatrix},\quad
J=\begin{pmatrix}1&1\\1&1\end{pmatrix},\quad c=2^t-1.
\]
The principal submatrix of
$\Phi=cI-E^{\otimes t}+F^{\otimes t}-J^{\otimes t}$
on the graph vertices is $L(G_n)$: the diagonal entry at $S$ is
$2^t-2-2^{t-|S|}$, and the off-diagonal entry is $-1$ for intersecting
subsets and zero otherwise.

With $e_0,e_1$ denoting the absence and presence basis, set
$w=e_0\otimes e_1-e_1\otimes e_0$.
For every $2\times2$ matrix $T$, one has
$T^{\otimes2}w=(\det T)w$. A paired factor $w$ thus contributes
$2,-1,0$ under $E,F,J$, respectively. It vanishes at the empty and
full subsets. Each subspace used below is invariant for $\Phi$ and is
supported on graph vertices, so its restriction is also a restriction of $L$.

If $t=2r+1$, $r\geq1$, use
$w^{\otimes r}\otimes\operatorname{span}\{e_0,e_1\}$.
Putting $a=2^r$ and $s=(-1)^r$, the restricted matrix is
\[
\begin{pmatrix}c-2a+s&s\\s&c-a\end{pmatrix},
\]
with discriminant $(a-s)^2+4$. The integer $h=a-s$ is at least $3$,
so $h^2<h^2+4<(h+1)^2$. Both eigenvalues are irrational.

If $t=2r$, $r\geq2$, use $r-1$ paired factors and the symmetric subspace
on the last two coordinates, with basis
$e_0e_0,(e_0e_1+e_1e_0)/\sqrt2,e_1e_1$.
Put $a=2^{r-1}$ and $s=(-1)^{r-1}$. The restriction is $cI-K$, where
\[
K=\begin{pmatrix}
4a-s&-s\sqrt2&-s\\-s\sqrt2&2a-s&0\\-s&0&a
\end{pmatrix}.
\]
The leading principal minors of $K-(a-1)I$ are
\[
3a+1-s,\quad(3a+1-s)(a+1-s)-2,\quad(3a-s)(a+1-s)-2.
\]
They are positive for $a\geq2$ and $s=\pm1$: the latter two are at
least $10$ and $8$, respectively. Sylvester's criterion gives
$\lambda_{\min}(K)>a-1$.
The principal submatrix of $K-aI$ on coordinates $1,3$ is
$\left(\begin{smallmatrix}3a-s&-s\\-s&0\end{smallmatrix}\right)$,
with determinant $-1$. A vector supported on those coordinates has a
negative Rayleigh quotient, giving $\lambda_{\min}(K)<a$.
Hence $c-\lambda_{\min}(K)$ is an eigenvalue of $L$ strictly between
the consecutive integers $c-a$ and $c-a+1$.
\end{proof}

\section{The family \texorpdfstring{$pq r^k$}{pq r to the k}}

\begin{theorem}
\label{thm:pqrk}
For distinct primes $p,q,r$ and every $k\geq1$, the graph $G_{pq r^k}$
has the irrational eigenvalues
\[
\frac{6k+1\pm\sqrt{(2k+2)^2-3}}2.
\]
In particular, every permutation of the exponent vector $(1,1,k)$
gives a nonintegral Laplacian spectrum.
\end{theorem}
\begin{proof}
The support classes $1,2$ have size $1$, and $13,23$ have size $k$.
Consider functions with values $s,-s$ on the first two classes,
$t,-t$ on the latter two, and zero elsewhere.
A vertex of support $1$ has $2k$ neighbors and Laplacian value
$2ks-kt$. A vertex of support $13$ has degree $4k$; after subtracting
the same-class and opposite-class values its Laplacian value is
$-s+(4k+1)t$. The swapped classes give their negatives, and all other
classes give zero by cancellation. The restricted matrix of the full
graph is therefore
\[
B_k=\begin{pmatrix}2k&-k\\-1&4k+1\end{pmatrix}.
\]
Its characteristic polynomial is $x^2-(6k+1)x+8k^2+k$, with
discriminant $(2k+2)^2-3$. Since
\[
(2k+1)^2<(2k+2)^2-3<(2k+2)^2
\]
for $k\geq1$, the two roots are irrational. Relabeling prime coordinates
preserves the graph.
\end{proof}

\section{A boundary family with integral antisymmetric block}

The following argument was supplied by Reza Nikandish in the public
collaboration discussion. We include it because it is the boundary case
needed for the linear cutoff in Section~\ref{sec:bounds}.

\begin{theorem}
\label{thm:boundary}
For every $a\geq2$, put $C=(a-1)(2a-1)$. For distinct primes $p,q,r$,
the graph $G_{p^a q^a r^C}$ is not Laplacian integral, although the
antisymmetric block~\eqref{eq:antisymmetric} has integer eigenvalues.
\end{theorem}
\begin{proof}
Specialize $f_{a,b}$ from~\eqref{eq:cubic} to $b=C$.
For $a=2$ it is $x^3-47x^2+702x-3312$. Its values at
$0,1,2,3,4$ modulo $5$ are $3,4,2,3,3$. Thus it has no root in
$\mathbb F_5$ and is irreducible over $\mathbb Q$; its real roots
are noninteger. Lemma~\ref{lem:join} gives noninteger graph eigenvalues.

For $a\geq3$ put $m=2a^3-2a^2+2a-1$. Expansion gives
\begin{align*}
f_{a,C}(m)&=2(5a^4-10a^3+7a^2-1),\\
f_{a,C}(m-1)&=-2(a-1)(a+2)(2a^4-7a^3+7a^2-a-3).
\end{align*}
The first value is positive because its parenthesis is
$5a^3(a-2)+7a^2-1$. Writing $z=a-3\geq0$, the last factor in the
second value is
$2z^4+17z^3+52z^2+68z+30>0$.
Consequently $f_{a,C}(m-1)<0<f_{a,C}(m)$, so a real root lies strictly
between $m-1$ and $m$. Its integer shift by $u=a^2C-1$ is noninteger.

The characteristic discriminant of $B_-$ is
$D=(a+1)^2b^2+2a(3a+1)b+a^2$.
At $b=C$ it equals $(2a^3-a^2+a-1)^2$, and the two eigenvalues are
$4a^3-3a^2+a$ and $2a^3-2a^2+1$. Thus the obstruction in this family
comes from the symmetric cubic.
\end{proof}

\section{Finite reduction and a linear cutoff}
\label{sec:bounds}

The repeated-exponent decomposition leaves two possible obstructions:
a nonsquare quadratic discriminant and a noninteger root of the symmetric
cubic. Exact factor pairs reduce the first exceptional set, while sign
intervals settle the large-exponent tail.

Fix $a,b\geq1$, and put
\[
h=a+1,\quad C=(a-1)(2a-1),\quad M=a^3(2a+1),\quad
A=h^2b+a(3a+1).
\]
The characteristic discriminant of~\eqref{eq:antisymmetric} is
\begin{equation}
\label{eq:discriminant}
D=h^2b^2+2a(3a+1)b+a^2.
\end{equation}
If $D$ is not a square, this block has irrational eigenvalues, which
remain irrational after the integer shift $a^2b-1$.

\begin{proposition}[Exact divisor reduction]
\label{prop:divisors}
For fixed $a\geq1$, the positive integers $b$ for which $D$ is a square
are exactly those obtained from positive factor pairs $d\leq e$ of $M$
such that
\begin{equation}
\label{eq:divisor-conditions}
b=\frac{d+e-a(3a+1)}{h^2}\in\mathbb Z_{\geq1},\qquad
s=\frac{e-d}{h}\in\mathbb Z_{\geq0}.
\end{equation}
In that case $D=s^2$.
\end{proposition}
\begin{proof}
The polynomial identity
\begin{equation}
\label{eq:product}
A^2-h^2D=4M
\end{equation}
holds for arbitrary $a,b$. If $D=s^2$ with $s\geq0$, then
$A-hs$ and $A+hs$ are positive: their product is positive and their sum
is $2A>0$. They have the same parity; since their product is divisible
by $4$, both are even. Write $A-hs=2d$ and $A+hs=2e$.
Then $de=M$, $d+e=A$ and $e-d=hs$, giving~\eqref{eq:divisor-conditions}.
Conversely those conditions yield $A=d+e$ and $hs=e-d$.
Using $(d+e)^2-(e-d)^2=4de=4M$ in~\eqref{eq:product} gives $D=s^2$.
\end{proof}

The divisor set is finite for every fixed $a$; this criterion has no
truncation in $b$. It gives first a quadratic bound.

\begin{corollary}
\label{cor:quadratic}
If $D$ is a square then $b\leq C$. Consequently every $b>C$ gives
a nonintegral Laplacian spectrum. For $a=1$, every $b\geq1$ is covered.
\end{corollary}
\begin{proof}
Since $(d-1)(e-1)\geq0$, $A=d+e\leq M+1$. The identity
\[
M+1-a(3a+1)=h^2C
\]
then gives $b\leq C$. The nonsquare case is covered by the antisymmetric
block. At $a=1$, $C=0$.
\end{proof}

\begin{lemma}
\label{lem:congruence}
Every admissible factor pair in Proposition~\ref{prop:divisors} satisfies
$d\equiv e\equiv1\pmod h$.
\end{lemma}
\begin{proof}
From $e-d=hs$ and $d+e=A$, reduction modulo $h$ gives
$e\equiv d$ and $2d\equiv2$, since $a(3a+1)\equiv2\pmod h$.
Also
\begin{equation}
\label{eq:strong-product}
(d-1)(e-1)=M-A+1=h^2(C-b).
\end{equation}
If $h$ is odd, the congruence $2(d-1)\equiv0$ gives the assertion.
If $h$ is even, the alternative to $d\equiv1$ is
$d-1\equiv h/2\pmod h$. Then $d-1$ and $e-1$ would both equal
$h/2$ times an odd integer. Their product divided by $h^2$ would be an
odd integer divided by $4$, contradicting~\eqref{eq:strong-product}.
This excludes the alternative for either parity of $a$.
\end{proof}

\begin{theorem}[Unified linear cutoff]
\label{thm:linear}
Let $a\geq2$ and
\[
R(a)=\frac{2(a-1)(a-2)}{a+2}=2a-10+\frac{24}{a+2}.
\]
For every integer $b>R(a)$ and distinct primes $p,q,r$, the graph
$G_{p^a q^a r^b}$ is not Laplacian integral. Every unresolved
repeated-exponent pair therefore has square $D$ and
$1\leq b\leq\lfloor R(a)\rfloor$.
\end{theorem}
\begin{proof}
If $D$ is nonsquare the antisymmetric block suffices. Otherwise use
Proposition~\ref{prop:divisors}. When $d=1$, $e=M$ and $b=C$;
Theorem~\ref{thm:boundary} settles this case.
If $d>1$, Lemma~\ref{lem:congruence} gives $d\geq T=a+2$.
Since $e\geq d\geq T$,
\[
\frac MT+T-(d+e)=(d-T)\left(\frac eT-1\right)\geq0.
\]
Hence
\[
b\leq\frac{M/T+T-a(3a+1)}{h^2}
=\frac{2(a-1)(a-2)}{a+2}=R(a).
\]
Thus $b>R(a)$ leaves either a nonsquare discriminant or the already
settled boundary case.
\end{proof}

At $(a,b)=(4,2)$ the square-discriminant pair has $d=6=a+2$,
so equality in the factor estimate is attained. The graph is still
nonintegral, as shown in the next section.

\subsection{Settling the first nonboundary factor}

The first possible value after $d=1$ is $d=a+2$.
Since $d$ divides $M$, the identity
\[
M=(a+2)(2a^3-3a^2+6a-12)+24
\]
forces $a+2\mid24$. With $a\geq2$ this leaves
$a=2,4,6,10,22$. Substitution into~\eqref{eq:divisor-conditions}
gives $b=0,2,5,12,35$, respectively. Excluding $b=0$, the complete
positive-$b$ list and its cubic certificates are:
\begin{center}
\begin{tabular}{clc}
\toprule
$(a,b)$ & $f_{a,b}(x)$ & Root-free prime\\
\midrule
$(4,2)$ & $x^3-86x^2+2304x-18816$ & 5\\
$(6,5)$ & $x^3-245x^2+19266x-488160$ & 7\\
$(10,12)$ & $x^3-842x^2+230160x-20534400$ & 13\\
$(22,35)$ & $x^3-4919x^2+7887858x-4132479120$ & 17\\
\bottomrule
\end{tabular}
\end{center}
Each polynomial has no root modulo its indicated prime. Its monic cubic
is therefore irreducible over $\mathbb Q$, and its real nonzero roots
lift to noninteger graph eigenvalues. Thus every possible $d=a+2$ pair
is settled, over all $a$, by this finite divisor argument.

\begin{theorem}[Second linear cutoff]
\label{thm:second}
For $a\geq2$ put
\[
S(a)=\frac{(a-3)(2a-3)}{2a+3}=a-6+\frac{27}{2a+3}.
\]
For every positive integer $b>S(a)$ and distinct primes $p,q,r$,
$G_{p^a q^a r^b}$ is not Laplacian integral.
Every remaining square-discriminant case has $d\geq2a+3$ and
$b\leq\lfloor S(a)\rfloor$.
\end{theorem}
\begin{proof}
Nonsquare $D$ is covered by the antisymmetric block. If $D$ is square,
the cases $d=1$ and $d=a+2$ are already settled. Every other admissible
factor satisfies $d\geq T=2a+3$ by Lemma~\ref{lem:congruence}.
Since $e\geq d\geq T$, the same factor inequality gives
\[
b\leq\frac{M/T+T-a(3a+1)}{(a+1)^2}
=\frac{(a-3)(2a-3)}{2a+3}=S(a).
\]
This proves the claim for $b>S(a)$.
\end{proof}

The new bound is strictly smaller than $R(a)$ for all $a\geq2$, since
\[
R(a)-S(a)=\frac{2a^3-a^2-a-6}{(a+2)(2a+3)}>0.
\]
Writing $z=a-2\geq0$, its numerator is $2z^3+11z^2+19z+4$.

\subsection{A finite reduction for each factor index}

For a nonboundary square-discriminant pair, write $d=1+hj$, $j\geq1$.
The following criterion fixes $j$ instead of $a$.

\begin{proposition}
\label{prop:index}
Fix $j\geq1$ and set $K_j=(j+1)^3(j+2)$.
Every admissible square-discriminant pair of index $j$ is obtained from
a positive divisor $d$ of $K_j$ with
$a=(d-1)/j-1\in\mathbb Z_{\geq2}$.
Set $e=M/d$, $v=(e-1)/(a+1)$, and retain exactly those candidates with
$d\leq e$ and $b=C-jv\geq1$. Then $e,v,b$ are integers and
$D=(v-j)^2$. Conversely every retained candidate is admissible.
\end{proposition}
\begin{proof}
Modulo $d=j(a+1)+1$ we have $ja\equiv-(j+1)$, whence
\[
j^4M=(ja)^3j(2a+1)\equiv(j+1)^3(j+2)=K_j\pmod d.
\]
Since $\gcd(d,j)=1$, we have $d\mid M$ if and only if $d\mid K_j$.
Thus $e$ is integral. Since $M\equiv d\equiv1\pmod h$, also
$e\equiv1\pmod h$, so $v$ is integral. The identity
\[
M=1+h(3a-2)+h^2C=(1+hj)(1+hv)
\]
gives $j+v+h jv=3a-2+hC$. Defining $b=C-jv$, we obtain
$d+e=h^2b+a(3a+1)$ and $e-d=h(v-j)$.
The ordering condition gives $v-j\geq0$; Proposition~\ref{prop:divisors}
now proves admissibility and $D=(v-j)^2$. The reverse direction follows
by applying these identities to any admissible index-$j$ pair.
\end{proof}

This is a complete finite problem for each fixed $j$, without a
truncation in $a$. Section~\ref{sec:completion} combines the first
three such problems with a uniform argument for all later indices.
For $j=2$, $K_2=108$ and $d=2a+3$ is an odd divisor of $108$ with
$d\geq7$. The only candidates are $d=9,27$, giving $a=3,12$.
The first has $e=21,v=5,b=0$ and is excluded. The second has
$e=1600,v=123,b=7$ and $\sqrt D=121$.
Thus $(12,7)$ is the only positive-$b$ case with $j=2$.

Its symmetric cubic is $f_{12,7}(x)=x^3-751x^2+180348x-13829760$.
Its values at all residues modulo $13$ are
$4,7,9,3,8,4,10,6,11,5,7,10,7$.
As a monic cubic with no root modulo $13$, it is irreducible over
$\mathbb Q$. Its real nonzero roots lift to noninteger graph eigenvalues,
settling the entire index-$2$ case.

\begin{theorem}[Third linear cutoff]
\label{thm:third}
For every $a\geq2$, put $T(a)=2(a-4)(a-2)/(3a+4)$.
If $b\geq1$ and $b>T(a)$, then $G_{p^a q^a r^b}$ is not Laplacian
integral. Every unresolved square-discriminant pair has $j\geq3$,
$d\geq3a+4$, and $b\leq\lfloor T(a)\rfloor$.
\end{theorem}
\begin{proof}
Nonsquare $D$ gives irrational antisymmetric eigenvalues.
For square $D$, the cases $j=0,1,2$ are all settled above.
Every other pair has $d\geq Q=3a+4$, and $e\geq d\geq Q$ gives
\[
b\leq\frac{M/Q+Q-a(3a+1)}{(a+1)^2}
=\frac{2(a-4)(a-2)}{3a+4}=T(a).
\]
This proves the claimed implication.
\end{proof}

For $a\geq3$ the new bound is strictly below $S(a)$, since
\[
S(a)-T(a)=\frac{2a^3-a^2-5a-12}{(2a+3)(3a+4)}>0.
\]
With $z=a-3\geq0$, the numerator is $2z^3+17z^2+43z+18$.
The exact expression
$T(a)=2a/3-44/9+320/(9(3a+4))$ shows its leading term.

\section{Completion of the repeated-exponent family}
\label{sec:completion}

The preceding factor-index reduction makes the repeated-exponent problem
finite at its first few levels. A uniform cubic sign interval handles every
later level, completing the family without a bound on either exponent.

\begin{theorem}
\label{thm:repeated}
For distinct primes $p,q,r$ and all integers $a,b\geq1$, the graph
$G_{p^a q^a r^b}$ is not Laplacian integral. The same conclusion holds
for every permutation of the exponent vector $(a,a,b)$.
\end{theorem}

The arithmetic reduction and a uniform sign obstruction for the cubic
combine to exhaust the family. We first settle one last finite factor level.

\begin{lemma}
\label{lem:third-index}
Every square-discriminant pair with index $j=3$ gives a nonintegral
Laplacian spectrum.
\end{lemma}
\begin{proof}
Proposition~\ref{prop:index} gives $d\mid K_3=320$ and $d=3a+4$.
For $a\geq2$, the complete divisor list of this shape is
$10,16,40,64,160$. The resulting values are
\begin{center}
\begin{tabular}{rrrrrl}
\toprule
$d$ & $a$ & $e$ & $v$ & $b$ & Outcome\\
\midrule
10 & 2 & 4 & 1 & 0 & $d>e$\\
16 & 4 & 36 & 7 & 0 & $b=0$\\
40 & 12 & 1080 & 83 & 4 & Root-free modulo $5$\\
64 & 20 & 5125 & 244 & 9 & Root-free modulo $7$\\
160 & 52 & 92274 & 1741 & 30 & Root-free modulo $7$\\
\bottomrule
\end{tabular}
\end{center}
The three positive-$b$ cubics are
\begin{align*}
f_{12,4}(x)&=x^3-568x^2+100032x-5248512,\\
f_{20,9}(x)&=x^3-1769x^2+992820x-174744000,\\
f_{52,30}(x)&=x^3-13394x^2+57771168x-80229951360.
\end{align*}
Their values at all residues of the indicated primes are, respectively,
$(3,3,3,4,2)$, $(3,2,4,1,6,4,1)$ and $(1,6,4,1,3,2,4)$.
Each monic cubic is therefore irreducible over $\mathbb Q$ and has
noninteger real roots. Section~\ref{sec:graph} supplies the lift to
noninteger graph eigenvalues.
\end{proof}

\begin{lemma}[A uniform consecutive-integer interval]
\label{lem:sign}
For $a,b\geq1$, put
\[
L(a)=\frac{a^3+2a^2+2a+1}{a(2a+1)}.
\]
If $b<L(a)$, then $f_{a,b}$ has a root strictly between $2ab-1$ and $2ab$.
\end{lemma}
\begin{proof}
Substitution in~\eqref{eq:cubic} gives
\begin{align*}
f_{a,b}(2ab)&=2a^2b^2>0,\\
f_{a,b}(2ab-1)&=-(a+b+1)
\bigl(a^3+2a^2+2a+1-a(2a+1)b\bigr).
\end{align*}
The second value is strictly negative when $b<L(a)$. The intermediate
value theorem gives the asserted noninteger, nonzero root.
\end{proof}

\begin{proof}[Proof of Theorem~\ref{thm:repeated}]
For $a=1$, Corollary~\ref{cor:quadratic} settles every positive $b$.
For $a\geq2$, nonsquare $D$ gives irrational antisymmetric eigenvalues.
When $D$ is square, index $j=0$ is Theorem~\ref{thm:boundary},
and indices $j=1,2$ are settled in Section~\ref{sec:bounds}.
Lemma~\ref{lem:third-index} settles $j=3$.

It remains to consider $j\geq4$. Then $d\geq Q=4a+5$, and
$e\geq d\geq Q$ implies
\[
\frac MQ+Q-(d+e)=(d-Q)\left(\frac eQ-1\right)\geq0.
\]
Since $d+e=(a+1)^2b+a(3a+1)$, we obtain
\[
b\leq B_4(a)=\frac{M/Q+Q-a(3a+1)}{(a+1)^2}
=\frac{(a-5)(2a-5)}{4a+5}.
\]
For every $a\geq2$,
\[
L(a)-B_4(a)=
\frac{41a^3-17a^2-11a+5}{a(2a+1)(4a+5)}>0.
\]
Indeed, with $z=a-2\geq0$, the numerator becomes
$41z^3+229z^2+413z+243$. Thus $b\leq B_4(a)<L(a)$, and
Lemma~\ref{lem:sign} gives a noninteger cubic root, which lifts across
the universal vertices. These cases exhaust all positive $a,b$.
Permuting the prime factors preserves the graph's isomorphism type.
\end{proof}

The preceding cutoffs and fixed-exponent theorems are intermediate
consequences of the arithmetic reduction. The uniform interval now
settles the infinite range left by those cutoffs, without a parameter scan.

\section{A unit exponent}
\label{sec:unit}

The complement quotient replaces coordinate symmetry when the other
two exponents are arbitrary. A root between zero and one then settles the
entire unit-exponent family after the small exceptions are identified.

\begin{theorem}
\label{thm:unit}
For distinct primes $p,q,r$ and every $b,c\geq1$, the graph
$G_{p q^b r^c}$ is not Laplacian integral. The same holds for
permutations of the exponent vector $(1,b,c)$.
\end{theorem}

The equal-exponent decomposition is replaced by a complement identity
on all six support classes. For arbitrary $a,b,c\geq1$, remove the
$u=abc-1$ universal vertices and denote the remaining graph by $H$.
In the order $1,2,3,12,13,23$, its class-size vector is
$w=(a,b,c,ab,ac,bc)^T$ and its vertex count is
$N=a+b+c+ab+ac+bc$. Disjoint supports are adjacent in the complement
$\overline H$, whose cell-constant Laplacian quotient is
\[
C=\begin{pmatrix}
b+c+bc&-b&-c&0&0&-bc\\
-a&a+c+ac&-c&0&-ac&0\\
-a&-b&a+b+ab&-ab&0&0\\
0&0&-c&c&0&0\\
0&-b&0&0&b&0\\
-a&0&0&0&0&a
\end{pmatrix}.
\]
Multiplication on the left by $\operatorname{diag}(w)$ makes $C$
symmetric. Moreover $C\mathbf1=0$ and $w^TC=0$.
If $B$ is the quotient of $H$, the complement relation gives
\begin{equation}
\label{eq:complement}
B+C=NI-\mathbf1w^T.
\end{equation}
Thus an eigenvector $Cv=\mu v$ with $\mu\ne0$ satisfies $w^Tv=0$
and $Bv=(N-\mu)v$. When $N-\mu\ne0$, Lemma~\ref{lem:join}
lifts it to an eigenvalue $N+u-\mu$ of $G_n$.

\begin{proof}[Proof of Theorem~\ref{thm:unit}]
Write $\det(xI-C)=x h_{a,b,c}(x)$. Direct determinant expansion yields
\begin{align*}
h_{a,b,c}(0)&=-abc(a+b+c)N,\\
h_{a,b,c}(a)&=-a^2b^2c^2(a^2+2a-b-c).
\end{align*}
At $a=1$, the first value is strictly negative and
\[
h_{1,b,c}(1)=b^2c^2(b+c-3).
\]
When $b+c>3$, a real root $\mu$ lies strictly between $0$ and $1$.
Equation~\eqref{eq:complement} and the universal-vertex lift produce a
Laplacian eigenvalue $V-\mu\in(V-1,V)$, where $V=N+u$ is the integer
vertex count. This eigenvalue is noninteger. When $b+c\leq3$, the
positive pairs are $(1,1),(1,2),(2,1)$; Theorem~\ref{thm:pqrk}
settles the corresponding permutations of $(1,1,k)$. Permuting the
prime factors preserves the graph's isomorphism type.
\end{proof}

The unit-exponent obstruction extends to any number of prime factors
by a variational argument on the complement graph.

\begin{theorem}
\label{thm:unit-any-prime-count}
Let $n=\prod_{i=1}^t p_i^{k_i}$ with $t\geq3$, distinct primes
$p_i$ and positive integer exponents $k_i$. If some $k_i=1$, then
$G_n$ has a Laplacian eigenvalue in $(V-1,V)$, where
$V=\prod_{i=1}^t(k_i+1)-2$. In particular, $G_n$ is not
Laplacian integral.
\end{theorem}
\begin{proof}
Permute the factors so that $k_t=1$, and put
\[
 A=\prod_{i<t}k_i,\qquad
 B=\prod_{i<t}(k_i+1)-1-A.
\]
The quantity $B$ is the sum of the weights of all nonempty proper
supports on the first $t-1$ coordinates. Since $t-1\geq2$, it is
positive. Remove the $A-1$ full-support universal vertices and let
$F=\overline H$. Then $N=|H|=A+2B+1$ and $V=N+A-1=2A+2B$.
The graph $F$ is connected: its singleton support classes form a
connected complete multipartite graph, and every proper support has
a singleton outside it.

The support $\{t\}$ consists of a single vertex $z$. Each of the
$A$ vertices supported on $[t-1]$ has $z$ as its sole neighbor in
$F$. The remaining $2B$ vertices form two collections of weight $B$:
nonempty proper supports $S\subset[t-1]$, and their counterparts
$S\cup\{t\}$. The vertex $z$ is adjacent to all vertices of the
first collection and none of the second.

Give a function $f$ the value $2B$ on the $A$ leaves, zero at $z$,
and $-A$ on the other $2B$ vertices. It is nonzero and
orthogonal to the constant vector. Its Laplacian energy and norm are
\[
 f^TL(F)f=A(2B)^2+BA^2=AB(4B+A),\qquad
 f^Tf=A(2B)^2+2BA^2=2AB(2B+A).
\]
Indeed, only edges incident to $z$ have unequal endpoint values.
Connectedness and the variational characterization therefore give
\[
 0<\lambda_2(F)\leq\frac{4B+A}{4B+2A}<1.
\]
For a corresponding eigenvector, the complement identity
$L(H)+L(F)=NI-J$ gives the nonzero eigenvalue $N-\lambda_2(F)$
of $H$. Lemma~\ref{lem:join} lifts it to
$V-\lambda_2(F)\in(V-1,V)$, proving the claim.
\end{proof}

The restriction $t\geq3$ is sharp for this conclusion. If $t=2$
and $k_2=1$, then $B=0$ and $F=K_{1,A}$, so the complement and
universal-vertex identities give an integral spectrum for $G_n$.
Theorem~\ref{thm:unit-any-prime-count} also gives a proof of
Theorems~\ref{thm:squarefree} for $t\geq3$ and~\ref{thm:unit}
without their separate spectral decompositions or small exceptions.

The general six-support reduction also gives an exact diagnostic for
fully distinct triples. For $1\leq a<b<c\leq7$, all thirty-five quintics
have a nonlinear irreducible factor over $\mathbb Q$. The quintics
at $(1,2,7)$ and $(2,3,5)$ have factor degrees $(1,4)$, and that at
$(2,4,6)$ has factor degrees $(2,3)$;
irreducibility of the entire quintic is therefore not a
uniform obstruction. Theorem~\ref{thm:unit} covers every triple beginning
with $1$ without a range restriction; the thirty-five-case calculation
itself remains finite evidence.

\section{A minimum exponent of two}
\label{sec:minimum-two}

Specializing the complement quintic to a minimum exponent of two
reduces its endpoint signs to positive polynomial expansions. One remaining
pair is handled by a rational interval, completing this minimum uniformly.

\begin{theorem}
\label{thm:minimum-two}
For distinct primes $p,q,r$ and all $b,c\geq1$, the graph
$G_{p^2q^br^c}$ is not Laplacian integral. Consequently every
three-prime exponent vector with minimum entry at most two gives a
nonintegral graph.
\end{theorem}

\begin{proof}
Use the complement quotient and quintic from Section~\ref{sec:unit}.
For $a=2$, $h_{2,b,c}(0)<0$. Putting $s=b+c$ and $t=bc$ gives
\[
h_{2,b,c}(1)=(s-13)t^2-(2s+6)t+3s^3+s^2-3s-1.
\]
By symmetry order $b\leq c$. A unit exponent is covered by
Theorem~\ref{thm:unit}; if $b=2$ or $b=c$, use
Theorem~\ref{thm:repeated}. It remains to consider $3\leq b<c$.

When $b=3$ and $c=5+v$, $v\geq0$, the exact expansion is
\[
h_{2,3,5+v}(1)=12v^3+112v^2+268v+120>0.
\]
When $b\geq4$ and $c>b$, write $b=4+u$ and $c=5+u+v$, with
$u,v\geq0$. Then
\begin{align*}
h_{2,4+u,5+u+v}(1)
={}&(u^2+8u+19)v^3\\
&+(4u^3+38u^2+128u+170)v^2\\
&+(5u^4+62u^3+271u^2+502u+368)v\\
&+2u^5+32u^4+190u^3+504u^2+552u+160>0.
\end{align*}
All coefficients are positive. In either region the opposite signs at
$0$ and $1$ give a root $\mu\in(0,1)$ and, by~\eqref{eq:complement}
and the universal-vertex lift, a noninteger graph eigenvalue in $(V-1,V)$.

The only remaining pair is $(b,c)=(3,4)$. Here
\[
h_{2,3,4}(x)=x^5-53x^4+953x^3-6523x^2+13134x-7560,
\]
and $h_{2,3,4}(1)=-48$ while
$h_{2,3,4}(3/2)=13437/32>0$. A root lies in $(1,3/2)$.
Since $V=58$, its lifted graph eigenvalue lies in $(113/2,57)$,
an interval containing no integer. This exhausts all positive $b,c$.
Combining with Theorem~\ref{thm:unit} proves the minimum-entry assertion.
\end{proof}

\section{A uniform cutoff and a minimum exponent of three}
\label{sec:distinct-tail}

\begin{proposition}
\label{prop:distinct-tail}
For distinct primes $p,q,r$ and positive exponents ordered as
$2\leq a\leq b\leq c$, the graph $G_{p^a q^b r^c}$ is not
Laplacian integral whenever $c\geq4a^2-2a$.
\end{proposition}

\begin{proof}
Use the complement quintic $h=h_{a,b,c}$ of Section~\ref{sec:unit}.
Its first endpoint is the cubic
\begin{equation}
\label{eq:general-first-endpoint}
h(1)=Lc^3+Qc^2+Rc+S,
\end{equation}
where
\begin{align*}
L={}&a^2+b^2-1,\\
Q={}&a^3-4a^2b^2+2a^2b-a^2+2ab^2+ab-3a
       +b^3-b^2-3b+3,\\
R={}&(a-1)(b-1)(2ab+3a+3b-3),\\
S={}&(a-1)(b-1)(a+b-1)(ab+a+b-1).
\end{align*}
Here $L,R,S>0$. In fact the sharper condition $Lc+Q\geq0$
already gives $h(1)>0$. Put $T(a)=4a^2-2a$. Then
\[
Q+T(a)L=b^2(b-1)+(2a^2+a-3)b+K(a),
\]
where $K(a)=4a^4-a^3-5a^2-a+3$. For $a=2+u$, $u\geq0$,
\[
K(2+u)=4u^4+31u^3+85u^2+95u+37>0.
\]
Consequently $Lc+Q>0$ for $c\geq T(a)$. Since $h(0)<0$,
a root lies in $(0,1)$, producing a noninteger graph eigenvalue in
$(V-1,V)$ by~\eqref{eq:complement} and Lemma~\ref{lem:join}.
\end{proof}

\begin{theorem}
\label{thm:minimum-three}
For distinct primes $p,q,r$ and all $b,c\geq1$, the graph
$G_{p^3q^br^c}$ is not Laplacian integral. Consequently every
three-prime exponent vector with minimum entry at most three gives a
nonintegral graph.
\end{theorem}

\begin{proof}
Theorems~\ref{thm:unit}, \ref{thm:minimum-two} and \ref{thm:repeated}
leave $4\leq b<c$. Proposition~\ref{prop:distinct-tail} settles
$c\geq30$. At $a=3$, equation~\eqref{eq:general-first-endpoint} is
\begin{align*}
h_{3,b,c}(1)={}&(b^2+8)c^3+(b-1)(b^2-30b-12)c^2\\
&+6(b-1)(3b+2)c+4(b-1)(b+2)(2b+1).
\end{align*}
The remaining domain $4\leq b<c<30$ has exactly
$\binom{26}{2}=325$ pairs. Exhaustive integer evaluation of this
cubic gives $257$ positive and $68$ negative values, with no zero.
Table~\ref{tab:minimum-three-signs} lists every negative range;
all other pairs in this finite domain have positive values.
The complete polynomial and separate Horner sign certificates are in
\path{results/distinct-tail.json}; their construction is described in
Section~\ref{sec:verification}. Thus a positive value gives a root in $(0,1)$.

\begin{table}[ht]
\centering
\begin{tabular}{rrrrrrrrrrr}
\toprule
$b$&4&5&6&7&8&9&10&11&12&13\\
\midrule
$c_{\min}$&5&6&7&8&9&10&11&12&13&14\\
$c_{\max}$&13&15&16&17&17&16&16&15&14&14\\
\bottomrule
\end{tabular}
\caption{Every negative $h_{3,b,c}(1)$ in $4\leq b<c<30$;
each column includes all integers between its two endpoints.}
\label{tab:minimum-three-signs}
\end{table}

For the negative values, the second endpoint is positive except at
$(b,c)=(4,5)$. Indeed, for $v\geq0$,
\[
h_{3,4,6+v}(2)=104v^3+1238v^2+4018v+2344>0.
\]
For $b\geq5$, write $b=5+u$, $c=6+u+v$, with $u,v\geq0$.
Direct expansion gives
\begin{align*}
h_{3,5+u,6+u+v}(2)
={}&(5u^2+54u+153)v^3\\
&+(20u^3+262u^2+1179u+1863)v^2\\
&+(25u^4+426u^3+2642u^2+7011u+6624)v\\
&+10u^5+218u^4+1828u^3+7200u^2+12600u+6480>0.
\end{align*}
These regions exhaust $4\leq b<c$ except $(4,5)$. Whenever
$h(1)<0<h(2)$, a root lies in $(1,2)$, giving a noninteger graph
eigenvalue in $(V-2,V-1)$.

Finally
\[
h_{3,4,5}(x)=x^5-83x^4+2327x^3-24133x^2+60576x-42480.
\]
Here $h(1)=-3792$ and $h(3/2)=48825/32>0$. The root in
$(1,3/2)$ lifts to a graph eigenvalue in $(233/2,117)$, since $V=118$.
Every interval used contains no integer. The finite nonvanishing
certificate and the written infinite-range arguments exhaust all $b,c$;
Theorems~\ref{thm:unit} and \ref{thm:minimum-two} give the final assertion.
\end{proof}

\section{An inertia criterion and four families with bounded exponent gaps}
\label{sec:low-spectrum}

An endpoint sign change need not detect two roots in the same interval.
For example, the complement quintic at equal exponents factors as
\[
h_{a,a,a}(x)=(x-a^2-a)\bigl(x^2-(a^2+4a)x+3a^2\bigr)^2.
\]
We instead count low eigenvalues by a symmetric Schur complement.

\begin{proposition}
\label{prop:low-spectrum}
Order the exponents as $a\leq b\leq c$, and let $m$ be an integer
with $2\leq m<a$. If
\begin{equation}
\label{eq:upper-inertia}
(c-m)(a+b+c-m)<mab,
\end{equation}
then the six-support complement quotient has exactly two nonzero
eigenvalues in $(0,m)$. If, in addition,
\begin{equation}
\label{eq:lower-inertia}
(a-m+1)(a+b+c-m+1)>(m-1)bc,
\end{equation}
both lie in $(m-1,m)$, so $G_{p^a q^b r^c}$ is not Laplacian integral.
\end{proposition}

\begin{proof}
Put $s=a+b+c$ and $P=abc$. Symmetrize $C$ by the square roots of its
cell sizes and reorder the classes as $1,2,3,23,13,12$. The resulting
real symmetric matrix is
\[
\mathcal C=\begin{pmatrix}A&-\sqrt P I\\-\sqrt P I&D\end{pmatrix},
\quad D=\operatorname{diag}(a,b,c),
\]
where $A$ has diagonal $(b+c+bc,a+c+ac,a+b+ab)$ and
off-diagonal entries $A_{ij}=-\sqrt{a_i a_j}$, with
$(a_1,a_2,a_3)=(a,b,c)$. For $0<x<a$, the block $D-xI$ is positive
definite, and Sylvester's law of inertia gives
\[
\operatorname{inertia}(\mathcal C-xI)
=\operatorname{inertia}(D-xI)+\operatorname{inertia}(K(x)),
\]
where
\[
K(x)=\operatorname{diag}(k_a(x),k_b(x),k_c(x))-vv^T,
\quad v=(\sqrt a,\sqrt b,\sqrt c)^T,
\quad k_t(x)=s-x-\frac{xP}{t(t-x)}.
\]
For fixed $x>0$, the function $k_t(x)$ increases with $t>x$, since
\[
\frac{\partial k_t(x)}{\partial t}
=\frac{xP(2t-x)}{t^2(t-x)^2}>0.
\]
Condition~\eqref{eq:upper-inertia} is exactly $k_c(m)<0$.
Thus all three diagonal values are negative, so $K(m)$ is negative
definite. The matrix $\mathcal C-mI$ has three negative and three
positive eigenvalues. The complement support graph is a triangle with
one pendant cell at each vertex; it is connected, so $C$ has a simple
zero eigenvalue and all other eigenvalues are positive. Exactly two
nonzero eigenvalues therefore lie in $(0,m)$.

For $x=m-1$, condition~\eqref{eq:lower-inertia} says $k_a(x)>0$,
so all three $k_t(x)$ are positive. Congruence by the inverse square
roots of these values turns $K(x)$ into $I-zz^T$, where
\[
\|z\|^2=\frac{a}{k_a(x)}+\frac{b}{k_b(x)}+\frac{c}{k_c(x)}>1,
\]
because each $k_t(x)<s$. Consequently $K(x)$ has one negative and two
positive eigenvalues, and the only eigenvalue of $C$ below $x$ is zero.
Both low positive eigenvalues lie in $(m-1,m)$. They are less than
$a<N$, so the complement identity and universal-vertex lift apply.
Their graph eigenvalues lie in $(V-m,V-m+1)$, which contains no integer.
\end{proof}

\begin{theorem}
\label{thm:bounded-gaps}
For every $a\geq1$, each of the exponent triples
\[
(a,a+1,a+2),\quad(a,a+1,a+3),\quad
(a,a+2,a+3),\quad(a,a+3,a+4)
\]
gives a nonintegral ideal intersection graph on three distinct primes.
\end{theorem}

\begin{proof}
For $a\leq3$, use Theorem~\ref{thm:minimum-three}. Assume $a\geq4$
and write the triple as $(a,a+d,a+e)$, where
$(d,e)=(1,2),(1,3),(2,3),(3,4)$. The left side minus the right side
of~\eqref{eq:upper-inertia}, at $m=3$, is respectively
\[
-6a,\quad-2a,\quad-4a,\quad4-2a.
\]
These values are negative. Proposition~\ref{prop:low-spectrum}
therefore places exactly two nonzero complement eigenvalues in $(0,3)$.
It remains to exclude the integers $1$ and $2$.

The first endpoints have the identities
$h_{a,a+d,a+e}(1)=-4M_{d,e}(a)P_{d,e}(a)$ from
Table~\ref{tab:gap-first-endpoints}. Every multiplier is positive for
$a\geq4$, and every coefficient of $P_{d,e}(4+u)$ is positive.
Thus $h(1)<0$ in each family.

\begin{table}[ht]
\centering
\begin{tabular}{lll}
\toprule
$(d,e)$&$M_{d,e}(a)$&Descending coefficients of $P_{d,e}(4+u)$\\
\midrule
$(1,2)$&$a(a+1)$&$(1,18,116,314,294)$\\
$(1,3)$&$a(a+1)(a+3)$&$(1,13,51,57)$\\
$(2,3)$&$a+2$&$(1,25,242,1126,2484,2022)$\\
$(3,4)$&$a+3$&$(1,28,302,1550,3699,3126)$\\
\bottomrule
\end{tabular}
\caption{Exact first-endpoint identities and positive shifted factors.}
\label{tab:gap-first-endpoints}
\end{table}

The second endpoints factor as
\begin{align*}
h_{a,a+1,a+2}(2)&=-aF_{1,2}(a),\\
h_{a,a+1,a+3}(2)&=-(a-2)(a+1)F_{1,3}(a),\\
h_{a,a+2,a+3}(2)&=-aF_{2,3}(a),\\
h_{a,a+3,a+4}(2)&=-F_{3,4}(a),
\end{align*}
where
\begin{align*}
F_{1,2}(a)&=a^5-9a^4-14a^3+93a^2+67a+6,\\
F_{1,3}(a)&=a^4-6a^3-42a^2-39a-10,\\
F_{2,3}(a)&=a^5-5a^4-48a^3-59a^2-39a-54,\\
F_{3,4}(a)&=a^6-a^5-74a^4-375a^3-967a^2-1220a-340.
\end{align*}
Table~\ref{tab:gap-modular} evaluates each factor at every residue of
one prime. No value is zero, so no factor has an integer root. The
linear multipliers are nonzero for $a\geq4$, hence $h(2)\ne0$.
Both positive eigenvalues in $(0,3)$ are therefore noninteger, and
their graph lifts are noninteger as well.

\begin{table}[ht]
\centering
\begin{tabular}{lll}
\toprule
Factor&Prime&Values at residues $0,\ldots,p-1$\\
\midrule
$F_{1,2}$&7&$(6,4,1,5,6,5,1)$\\
$F_{1,3}$&11&$(1,3,9,8,2,6,4,4,5,8,5)$\\
$F_{2,3}$&7&$(2,6,5,3,5,4,3)$\\
$F_{3,4}$&13&$(11,1,3,4,3,1,9,9,11,11,6,1,8)$\\
\bottomrule
\end{tabular}
\caption{Complete root-free modular certificates for the second endpoints.}
\label{tab:gap-modular}
\end{table}
\end{proof}

\begin{corollary}
\label{cor:span-three}
Every positive three-prime exponent triple with maximum minus minimum
at most three gives a nonintegral ideal intersection graph.
\end{corollary}

\begin{proof}
Repeated entries are covered by Theorem~\ref{thm:repeated}. After
sorting distinct entries, the only possibilities have gaps
$(d,e)=(1,2),(1,3),(2,3)$, all covered by
Theorem~\ref{thm:bounded-gaps}.
\end{proof}

\subsection{The complete finite region for minimum exponents four through seven}

The uniform cutoff makes the requested diagnostic range complete for each
of the minimum exponents $4,5,6,7$. Here the relevant quintic is
$g_B(x)=\det(xI-B)/x$ for $H$, whereas the preceding endpoint arguments
use the complement quintic $h_C(x)=-g_B(N-x)$.

\begin{corollary}
\label{cor:minimum-seven}
Every positive three-prime exponent triple with minimum entry at most
seven gives a nonintegral ideal intersection graph.
\end{corollary}

\begin{proof}
Theorem~\ref{thm:minimum-three} handles minimum entries at most three.
For a minimum entry $a\in\{4,5,6,7\}$, repeated entries are covered by
Theorem~\ref{thm:repeated}, and Proposition~\ref{prop:distinct-tail}
settles $c\geq4a^2-2a$. The remaining ordered triples lie in the exact
finite domain $a<b<c<4a^2-2a$.

For all $27\,562$ triples in this domain, the integer discriminant
$\Delta_{a,b,c}$ of $g_B$ is positive and strictly between two
consecutive integer squares. Table~\ref{tab:open-region} gives the
complete counts. The certificate
\path{results/open-region-discriminants.csv} records every triple,
discriminant and integer $q$ satisfying
$q^2<\Delta_{a,b,c}<(q+1)^2$. Each discriminant is independently
recomputed as a $9\times9$ integer Sylvester determinant by
\path{scripts/check_open_region_certificate.py}, using fraction-free
elimination and exact divisions. The certificate exhausts the domain
without omitted or duplicate triples.

If all roots of the monic quintic $g_B$ were integers, its discriminant,
the product of the squared pairwise root differences, would be an
integer square. Thus every certified quintic has a noninteger root.
The quotient is similar to a real symmetric Laplacian, so the root is
real and nonzero; Lemma~\ref{lem:join} gives a noninteger graph
eigenvalue. This complete finite certificate and the written cutoff
exhaust all triples with minimum entry at most seven.
\end{proof}

\begin{table}[ht]
\centering
\begin{tabular}{rrrrrr}
\toprule
Minimum $a$&4&5&6&7&Total\\
\midrule
Exclusive cutoff&56&90&132&182&---\\
Complete pairs&1275&3486&7750&15051&27562\\
Nonsquare discriminants&1275&3486&7750&15051&27562\\
\bottomrule
\end{tabular}
\caption{Exhaustive discriminant certificates in the region derived from
the uniform cutoff.}
\label{tab:open-region}
\end{table}

The signs of $(g_B(1),g_B(2),g_B(3))$ are $(-,-,-)$ throughout this
finite range. The complement signs instead have three patterns,
$(-,+,-),(-,+,+),(+,+,+)$, recorded separately in the diagnostic.
Different endpoint patterns do not rule out a uniform proof, and equal
signs can miss two roots in one interval. Proposition~\ref{prop:low-spectrum}
addresses that issue without extrapolating from these finite counts.

\section{A uniform low eigenvalue and a balanced region}
\label{sec:mixed-inertia}

The Schur complement of Proposition~\ref{prop:low-spectrum} also
counts roots when its diagonal has mixed signs. A comparison at three
gives a low eigenvalue for every positive real parameter triple.

\begin{lemma}
\label{lem:mixed-inertia}
For $0<x<a$, suppose all three $k_t(x)$ are nonzero. Let $q$ be the
number of negative diagonal values and put
\[
F(x)=1-\frac{a}{k_a(x)}-\frac{b}{k_b(x)}-\frac{c}{k_c(x)}.
\]
Then $n_-(K(x))=q+\mathbf1_{\{F(x)<0\}}$ and
$\operatorname{nullity}K(x)=\mathbf1_{\{F(x)=0\}}$.
The number of positive complement quotient eigenvalues strictly below
$x$ is $n_-(K(x))-1$.
\end{lemma}

\begin{proof}
The bordered matrix
\[
\begin{pmatrix}\operatorname{diag}(k_t(x))&v\\v^T&1\end{pmatrix}
\]
is congruent both to $K(x)\oplus[1]$ and to
$\operatorname{diag}(k_t(x))\oplus[F(x)]$. Sylvester's law of inertia
gives the formulas. The block $D-xI$ is positive definite, and $C$ has
a simple zero eigenvalue; subtracting this eigenvalue gives the stated count.
\end{proof}

A direct comparison at three avoids the reciprocal formula at zero
diagonal entries and supplies a uniform low positive root.

\begin{theorem}
\label{thm:uniform-low-root}
For every ordered real triple $0<a\leq b\leq c$, the six-support
complement quotient has a nonzero eigenvalue in $(0,3)$.
\end{theorem}

\begin{proof}
Use $s=a+b+c$, $P=abc$ and $v=(\sqrt a,\sqrt b,\sqrt c)^T$ from
Proposition~\ref{prop:low-spectrum}, and define
\[
L=\operatorname{diag}\left(s-\frac{3P}{a^2},
s-\frac{3P}{b^2},s-\frac{3P}{c^2}\right)-vv^T.
\]
Expansion of this diagonal-minus-rank-one determinant gives
\begin{align*}
\det L&=6\bigl(c(a-b)^2+b(a-c)^2+a(b-c)^2\bigr),\\
v^TLv&=-3P\left(\frac1a+\frac1b+\frac1c\right)<0.
\end{align*}
The determinant and Rayleigh value give exactly two negative eigenvalues
at unequal parameters. At equal parameters $L=-vv^T$; thus $L$ is
nonpositive on a two-dimensional subspace in both cases.

For $a>3$, the identity
\[
K(3)=L-\operatorname{diag}_{t\in\{a,b,c\}}
\left(3+\frac{9P}{t^2(t-3)}\right)
\]
subtracts a positive definite matrix. Since $D-3I>0$, this gives
at least two negative eigenvalues of $\mathcal C-3I$.

If $b\leq3$, restrict $\mathcal C-3I$ to the first two singleton
supports and their paired supports. Its blocks are
$\bigl[\begin{smallmatrix}A_0&-\sqrt P I_2\\-\sqrt P I_2&D_0\end{smallmatrix}\bigr]$
with $D_0\leq0$. On vectors $(z,tz)$, the quadratic form is at most
$z^T(A_0-2t\sqrt P I_2)z$, which is negative definite for large $t$.

It remains to consider $a\leq3<b$. Eliminate the positive paired
diagonals $b-3,c-3$. On singleton supports $2,3$, the vector
$(\sqrt b,\sqrt c)$ has quadratic value
\[
 q=(a-3)(b+c)-3P\left(\frac1{b-3}+\frac1{c-3}\right)<0.
\]
Compress the remaining matrix to singleton $1$, its paired support,
and this vector. The resulting matrix is
\[
 \begin{pmatrix}M&-\sqrt P&r\\-\sqrt P&a-3&0\\r&0&q\end{pmatrix},
 \quad M=bc+b+c-3>0,\quad r=-\sqrt a(b+c).
\]
Its determinant $M(a-3)q-Pq-(a-3)r^2$ is positive.
Its positive and negative diagonal entries therefore force exactly
two negative eigenvalues. This includes the singular boundary $a=3$.

In every case $\mathcal C-3I$ has at least two negative eigenvalues.
Since $\mathcal C$ is positive semidefinite with a simple zero,
at least one positive eigenvalue lies strictly below three.
\end{proof}

\begin{corollary}
\label{cor:balanced-interval}
For $4\leq a\leq b\leq c$, if
\begin{equation}
\label{eq:balanced-interval}
(a-2)(a+b+c-2)>2bc,
\end{equation}
then $C$ has an eigenvalue in $(2,3)$, and $G_{p^a q^b r^c}$ is not
Laplacian integral. Nonzero values of both $h_C(1)$ and $h_C(2)$ also
imply nonintegrality. Hence $h_C(1)h_C(2)=0$ is necessary for an integral graph.
\end{corollary}

\begin{proof}
Condition~\eqref{eq:balanced-interval} is exactly $k_a(2)>0$, so all
three diagonal values are positive. As in
Proposition~\ref{prop:low-spectrum}, congruence turns $K(2)$ into
$I-zz^T$, with $\|z\|^2=\sum_t t/k_t(2)>1$, since every $k_t(2)<s$.
It has one negative and two positive eigenvalues and no zero eigenvalue.
Consequently the only quotient eigenvalue at or below two is zero.
The positive root below three from Theorem~\ref{thm:uniform-low-root}
lies in $(2,3)$ and lifts to $(V-3,V-2)$.
For the endpoint assertion, a positive root in $(0,3)$ can be an integer
only if it equals one or two, giving the necessary condition.
\end{proof}

\begin{corollary}
\label{cor:balanced-span}
Let $a$ be the minimum exponent and $\rho$ the maximum minus minimum.
Every positive three-prime exponent triple with
\[
a>\rho+4+\sqrt3(\rho+2)
\]
gives a nonintegral ideal intersection graph. In particular,
$a\geq3\rho+8$ still suffices.
\end{corollary}

\begin{proof}
Write the ordered triple as $(a,a+d,a+\rho)$ with $0\leq d\leq\rho$.
The left side minus the right side of~\eqref{eq:balanced-interval} is
\[
a^2-2a\rho-8a-2\rho^2-4\rho+4+(\rho-d)(a+2\rho+2).
\]
The last term is nonnegative, while the first part factors as
\[
\bigl(a-\rho-4-\sqrt3(\rho+2)\bigr)
\bigl(a-\rho-4+\sqrt3(\rho+2)\bigr).
\]
Both factors are positive under the stated bound, which also gives
$a>4+2\sqrt3>4$. Corollary~\ref{cor:balanced-interval} applies.
The old linear bound is larger, since
$3\rho+8-[\rho+4+\sqrt3(\rho+2)]=(2-\sqrt3)(\rho+2)>0$.
\end{proof}

The new threshold is exact for guaranteeing the strict diagonal
criterion~\eqref{eq:balanced-interval} from the minimum and span alone:
at fixed $a,\rho$ its left side minus right side is minimized by
$b=c=a+\rho$, and it vanishes at the displayed threshold.
This optimality concerns that sufficient criterion, not the entire
spectral region where a root lies in $(2,3)$. For example,
$(35,45,45)$ and $(281,381,381)$ satisfy the new bound but not
$a\geq3\rho+8$.

For example, $(20,21,24)$, $(38,43,48)$ and $(308,350,408)$ satisfy
the span criterion and lie outside the four fixed-gap families of
Theorem~\ref{thm:bounded-gaps}. The argument allows an unbounded span
and uses no exponent scan or modular endpoint certificate. The necessary
endpoint-zero condition is not sufficient for integrality: those
triples may still have other noninteger quotient eigenvalues.

\section{Arithmetic obstructions from the complement quotient}
\label{sec:arithmetic}

The uniform low-root bound rules out common exponent divisors greater
than two. Reduction of the characteristic polynomial gives further
infinite congruence classes, even when another quotient root is an integer.

\begin{theorem}
\label{thm:common-divisor}
Every positive three-prime exponent triple with $\gcd(a,b,c)\geq3$
gives a nonintegral ideal intersection graph.
\end{theorem}

\begin{proof}
Put $d=\gcd(a,b,c)$. All entries of the six-support complement quotient
$C$ are integer combinations of the exponents and their pairwise
products, so $C/d$ is an integer matrix. Every rational eigenvalue of
an integer matrix is an integer, by the rational-root theorem for its
monic characteristic polynomial. Therefore every integer eigenvalue
of $C$ is a multiple of $d$.

Order $a\leq b\leq c$. If $a=3$, Theorem~\ref{thm:minimum-three}
applies. Otherwise $a\geq4$, and Theorem~\ref{thm:uniform-low-root}
supplies an eigenvalue $0<\mu<3$. No positive multiple of $d\geq3$
lies in this interval. Hence $\mu$ and its graph lift $V-\mu$ are
noninteger.
\end{proof}

In particular, if $m$ has exactly three distinct prime factors, then
$G_{m^d}$ is not Laplacian integral for every $d\geq3$.

\begin{theorem}
\label{thm:congruence}
A positive three-prime exponent triple gives a nonintegral ideal
intersection graph under either of the following conditions:
\begin{enumerate}
\item all three exponents are odd;
\item there is an odd prime $\ell$ and a nonzero residue $t$ such that
$a\equiv b\equiv c\equiv t\pmod\ell$ and $t^2+8t+4$ is a quadratic
nonresidue modulo $\ell$.
\end{enumerate}
In particular, common residues $1,3,4$ modulo $5$, or $1,2,4,5$ modulo
$7$, imply nonintegrality.
\end{theorem}

\begin{proof}
The monic integer quintic $h_C(x)=\det(xI-C)/x$ would split into
linear factors modulo every prime if all its roots were integers.
A nonlinear irreducible factor in one reduction therefore supplies a
noninteger real complement root and a noninteger graph lift.

If the exponents are odd, reduce the general quotient modulo two:
\[
h_C(x)\equiv x(x^2+x+1)^2\pmod2.
\]
The quadratic has no root in $\mathbb F_2$, proving the first assertion.
For a common residue $t$, the equal-exponent identity gives
\[
h_C(x)\equiv(x-t^2-t)
\bigl(x^2-(t^2+4t)x+3t^2\bigr)^2\pmod\ell.
\]
The quadratic discriminant is $t^2(t^2+8t+4)$. For odd $\ell$ and
nonzero $t$, the assumed nonresidue makes the quadratic irreducible.
Modulo five the discriminant factors for $t=1,3,4$ are $3,2,2$,
respectively; the squares are $0,1,4$. Modulo seven the corresponding
values for $t=1,2,4,5$ are $6,3,3,6$, outside the square residues.
\end{proof}

These conditions allow unequal exponents with arbitrarily large ratios
or differences. For example, $(11,16,21)$, $(13,18,23)$ and $(14,19,24)$
have gcd one and mixed parity, but meet the modulo-five criterion.
The other parity reductions are $x^5$ for three even exponents and
$x^2(x+1)^3$ for one or two odd exponents. Their complete splitting
does not establish integrality.

The arithmetic arguments do not assume that the smallest positive
quotient root is noninteger. The necessary endpoint-zero condition of
Corollary~\ref{cor:balanced-interval} remains insufficient: the already
settled boundary triple $(9,9,136)$ has $h_C(1)=0$ and other noninteger
quotient roots. The remaining three-prime problem has gcd at most two,
at least one even exponent, and lies outside the congruence classes
above as well as the earlier proved regions.

\subsection{A mixed-parity divisibility obstruction}
\label{sec:mixed-parity-congruence}

The distribution of hypothetical integer roots modulo four gives further
infinite classes with unequal exponent valuations.

\begin{lemma}
\label{lem:three-odd-roots}
Let $f\in\mathbb Z[x]$ be monic of degree five with
$f(x)\equiv x^2(x+1)^3\pmod2$. If all its roots are integers,
then $32\mid f(1)$ or $32\mid f(3)$.
\end{lemma}

\begin{proof}
The binary factorization forces precisely three odd roots and two even
roots, counted with multiplicity. At least two odd roots share a residue
$r\in\{1,3\}$ modulo four. In $f(r)$, their two factors are divisible
by four and the third odd-root factor is divisible by two. The even-root
factors are odd. Therefore $32\mid f(r)$, including when $f(r)=0$.
\end{proof}

\begin{theorem}
\label{thm:mixed-parity-congruence}
Let $a,b$ be the two odd exponents and $c$ the even exponent of a positive
three-prime exponent triple. Each row in the following table, including
exchange of $a,b$, implies that its ideal intersection graph is nonintegral.
\begin{center}
\begin{tabular}{ccc|cc}
\toprule
$a\bmod8$ & $b\bmod8$ & $c\bmod8$ & $h_C(1)\bmod32$ & $h_C(3)\bmod32$\\
\midrule
1&3&2&8&16\\
1&7&2&24&16\\
3&7&2&16&8\\
3&5&6&8&16\\
3&7&6&16&24\\
5&7&6&24&16\\
\bottomrule
\end{tabular}
\end{center}
\end{theorem}

\begin{proof}
For either mixed exponent parity,
$h_C(x)\equiv x^2(x+1)^3\pmod2$. Substituting
$(a,b,c)=(8u+r_1,8v+r_2,8w+r_3)$ in the general quotient quintic gives
the last two columns: every nonconstant coefficient in $u,v,w$ is
divisible by $32$. Both entries are nonzero modulo $32$, contradicting
Lemma~\ref{lem:three-odd-roots} if all quotient roots were integers.
The resulting noninteger real complement eigenvalue has a noninteger
graph lift.
\end{proof}

For example, $(10,17,19)$ and $(14,19,21)$ meet these criteria with
gcd one and unequal $2$-adic valuations. There is no minimum or ratio
bound in the argument. The weaker divisibilities $4\mid h_C(2t)$
and $8\mid h_C(2t+1)$ hold identically for both mixed exponent parities;
they alone do not supply the contradiction. The distribution of the
three odd roots modulo four is the additional constraint.

\section{A linear bound for the endpoint two}
\label{sec:endpoint-reduction}

The two possible low integer roots have different arithmetic constraints.
The endpoint two is confined to a linear region in the minimum exponent;
this gives a new nonintegral tail for even exponent triples.

\begin{theorem}
\label{thm:endpoint-two-bound}
Let $4\leq a\leq b\leq c$, and put
\[
R(a)=\frac{7a^2-2a+8}{a+2}=7a-16+\frac{40}{a+2}.
\]
If $h_C(2)=0$, then $c<R(a)$. In fact, $c\geq R(a)$ implies
$h_C(2)>0$.
\end{theorem}

\begin{proof}
Fix $a,c$ and expand the symmetric endpoint polynomial in $b$:
\[
h_C(2)=A(a,c)b^3+B(a,c)b^2+D(a,c)b+E(a,c),
\]
where
\begin{align*}
A(a,c)&=ac(a+c-1)+2a(a-1)+2c(c-1)-4,\\
B(a,c)&=c^2\bigl((a+2)c-(7a^2-2a+8)\bigr)\\
&\quad +(a^3+2a^2-6a-4)c+2(a-2)(a^2-2a-6),\\
D(a,c)&=(a-2)(c-2)\bigl(a^2c+a^2+ac^2+6ac+4a+c^2+4c-12\bigr),\\
E(a,c)&=2(a-2)(c-2)(a+c-2)(ac+a+c-2).
\end{align*}
For $a,c\geq4$, the coefficients $A,D,E$ are strictly positive.
Both $a^3+2a^2-6a-4$ and $a^2-2a-6$ are positive at $a=4$
and increasing on $[4,\infty)$. Thus $c\geq R(a)$ also gives $B>0$,
which proves the strict endpoint sign. Conversely, an endpoint zero
requires $B<0$, so $(a+2)c<7a^2-2a+8$.
\end{proof}

\begin{corollary}
\label{cor:even-linear-tail}
If $4\leq a\leq b\leq c$ are all even and $c\geq R(a)$,
then $G_{p^a q^b r^c}$ is not Laplacian integral.
For minimum exponent at least eight, $c\geq7a-12$ suffices.
\end{corollary}

\begin{proof}
Every entry of $C$ is even, so every integer eigenvalue of $C$ is even
by the rational-root theorem applied to $C/2$. In particular, one is
not an eigenvalue. Theorem~\ref{thm:endpoint-two-bound} excludes two,
and Theorem~\ref{thm:uniform-low-root} provides a noninteger root in
$(0,3)$. Equivalently, $h_C(0)<0<h_C(2)$ directly yields a root in
$(0,2)$, which cannot equal one. Its graph lift is noninteger.
For $a\geq8$, $40/(a+2)\leq4$ gives $R(a)\leq7a-12$.
\end{proof}

The triples $(8,10,44)$, $(10,14,58)$ and $(40,42,266)$ meet this
criterion with gcd two. Their largest exponents are below the earlier
quadratic cutoff and their spans are outside the earlier balanced region.
The already settled repeated triple $(10,10,12)$ has $h_C(2)=0$ and
obeys the strict bound; an integer endpoint does not establish integrality.

\begin{proposition}
\label{prop:endpoint-divisors}
For fixed $a,b\geq4$, any positive integer $c$ on an endpoint-zero
surface satisfies
\begin{align*}
h_C(1)=0&\ \Longrightarrow\ c\mid(a-1)(b-1)(a+b-1)(ab+a+b-1),\\
h_C(2)=0&\ \Longrightarrow\ c\mid2(a-2)(b-2)(a+b-2)(ab+a+b-2).
\end{align*}
\end{proposition}

\begin{proof}
Both endpoints are integer cubics in $c$. Their constant terms are
the displayed nonzero products. Reduce each endpoint equation modulo $c$.
\end{proof}

For a fixed pair this gives finite divisor candidates, which must still
satisfy the endpoint equation. Reza's requested diagnostic for middle
exponents through fifteen now has a complete independent certificate.

\begin{corollary}
\label{cor:middle-fifteen}
Every positive ordered three-prime exponent triple $a\leq b\leq c$
with $b\leq15$ gives a nonintegral ideal intersection graph.
This conclusion uses the complete finite endpoint certificate below.
\end{corollary}

\begin{proof}
Repeated exponents are covered by Theorem~\ref{thm:repeated}, and
minimum exponents one, two and three by Theorems~\ref{thm:unit},
\ref{thm:minimum-two} and~\ref{thm:minimum-three}, respectively.
It remains to consider $4\leq a<b\leq15$ and $c>b$.

For each of these $66$ pairs, both endpoint polynomials are integer
cubics in $c$ with the positive constant terms in
Proposition~\ref{prop:endpoint-divisors}. The
\href{https://github.com/the-omega-institute/ideal-intersection-laplacian/blob/10ff01f0625c7264b8bc5bb33b1a9d7dcb88a93c/results/endpoint-surfaces-verification.json}{complete endpoint certificate}
records all $132$
coefficient lists and exhausts their $8\,527$ positive divisors using
integer trial division. Independent integer Horner evaluation at every
one of the $7\,585$ divisors above $b$ gives no zero. Thus neither
endpoint vanishes for any integer $c>b$, with no imposed upper bound
on $c$. Theorem~\ref{thm:uniform-low-root} supplies a positive root
in $(0,3)$; it can equal neither one nor two, proving nonintegrality.
\end{proof}

The reflected quotient in this diagnostic uses
$N=|V(H)|=a+b+c+ab+ac+bc$, not the full graph vertex count.
The inclusive requested bound $b\leq15$ and exact rational linear-factor
extraction are cross-checked against the independent divisor certificate.
The empty lists settle this pair range; they do not prove either
endpoint surface globally empty.

The remaining all-even triples have gcd two,
$8\leq a<b<c<R(a)$ and $b\geq16$, outside the preceding proved regions.
The mixed-parity surface $h_C(1)=0$ still needs another obstruction.

\section{Complete fixed-exponent families and exact diagnostics}

\begin{theorem}
\label{thm:small}
For every $a\in\{2,3,4\}$, every $b\geq1$ and distinct primes $p,q,r$,
the graph $G_{p^a q^a r^b}$ is not Laplacian integral.
\end{theorem}
\begin{proof}
The cutoff values are $R(2)=0$, $R(3)=4/5$, and $R(4)=2$.
Theorem~\ref{thm:linear} settles all $b$ for $a=2,3$, and all $b>2$
for $a=4$. At $(a,b)=(4,1)$, $D=145$ is nonsquare.
At $(4,2)$ the symmetric cubic is
\[
f_{4,2}(x)=x^3-86x^2+2304x-18816.
\]
Its values at $0,1,2,3,4$ modulo $5$ are $4,3,1,4,3$, so it has
no root modulo $5$. A cubic over a field without a root is irreducible.
Since this cubic is monic, reduction modulo $5$ proves irreducibility over
$\mathbb Q$. Its real nonzero roots lift to noninteger graph eigenvalues.
\end{proof}

\begin{corollary}
\label{cor:five-six}
The same all-$b$ conclusion holds for $a=5,6$. Thus it holds for every
$a\in\{2,3,4,5,6\}$.
\end{corollary}
\begin{proof}
The second cutoff gives $S(5)=14/13$ and $S(6)=9/5$, leaving only
$b=1$ in each case. The discriminants are $221$ and $313$, respectively.
Since $14^2<221<15^2$ and $17^2<313<18^2$, both give irrational
antisymmetric eigenvalues. Every larger $b$ is covered by
Theorem~\ref{thm:second}.
\end{proof}

\begin{corollary}
\label{cor:seven-nine}
Every positive $b$ also gives a nonintegral spectrum for $a=7,8,9$.
Thus every repeated exponent $a\in\{2,3,4,5,6,7,8,9\}$ is settled.
\end{corollary}
\begin{proof}
The third cutoff gives $T(7)=6/5$, $T(8)=12/7$, $T(9)=70/31$.
Only four pairs remain under these bounds:
\[
(a,b)=(7,1),(8,1),(9,1),(9,2).
\]
Their discriminants are $421,545,685,1489$, respectively, between the
consecutive squares with lower square roots $20,23,26,38$.
Each therefore gives irrational antisymmetric eigenvalues, and
Theorem~\ref{thm:third} covers every larger $b$.
\end{proof}

For comparison, the complete divisor certificates before applying the
linear cutoff are given in Table~\ref{tab:divisors}. The constants factor
as $M=40=2^3\cdot5$, $189=3^3\cdot7$, and
$576=2^6\cdot3^2$. Testing every divisor $d\leq\sqrt M$ in
\eqref{eq:divisor-conditions} gives exactly the displayed admissible pairs.
The corresponding cubics in Table~\ref{tab:cubics} have no root modulo
the listed primes. These finite certificates supply a second proof using
the complete exceptional list. Three pairs are in the boundary family;
$(4,2)$ is the additional symmetric-block obstruction.

\begin{table}[ht]
\centering
\begin{tabular}{rrrrr}
\toprule
$a$ & $M$ & $(d,e)$ & $b$ & $\sqrt D$\\
\midrule
2 & 40 & $(1,40)$ & 3 & 13\\
3 & 189 & $(1,189)$ & 10 & 47\\
4 & 576 & $(6,96)$ & 2 & 18\\
4 & 576 & $(1,576)$ & 21 & 115\\
\bottomrule
\end{tabular}
\caption{Every admissible factor pair for $a=2,3,4$.}
\label{tab:divisors}
\end{table}

\begin{table}[ht]
\centering
\begin{tabular}{clc}
\toprule
$(a,b)$ & $f_{a,b}(x)$ & Witness prime\\
\midrule
$(2,3)$ & $x^3-47x^2+702x-3312$ & 5\\
$(3,10)$ & $x^3-187x^2+11220x-214200$ & 13\\
$(4,2)$ & $x^3-86x^2+2304x-18816$ & 5\\
$(4,21)$ & $x^3-485x^2+75492x-3721536$ & 5\\
\bottomrule
\end{tabular}
\caption{Root-free modular reductions for the exceptional cubics.}
\label{tab:cubics}
\end{table}

The separate diagnostic range $a=2,3,4$, $1\leq b\leq30$ contains
$90$ cubics, each with a noninteger root certified in an interval between
consecutive integers. The reducible cases in that range are $b=2,6$
for $a=2$, $b=3,12$ for $a=3$, and $b=4,20$ for $a=4$.
Each has a linear and an irreducible quadratic factor.
These observations do not require whole-cubic irreducibility.
Indeed, for every $a\geq1$,
\begin{align*}
f_{a,a}(x)&=(x-2a(a+1))
 \bigl(x^2-(5a^2+2a)x+6a^4+3a^3\bigr),\\
f_{a,a(a+1)}(x)&=(x-2a(a+1)^2)
 \bigl(x^2-(3a^3+4a^2+2a)x+2a^6+6a^5+7a^4+3a^3\bigr).
\end{align*}
Consequently no prime can make every diagnostic cubic irreducible: the
monic factors in the diagonal family retain positive degree modulo every
prime.

For the second family the quadratic discriminant is $a^2(a^4+4a+4)$.
For $a\geq3$ the radicand is between $a^4$ and $(a^2+1)^2$, since
$2a^2-4a-3=2(a-3)^2+8(a-3)+3>0$.
At $a=2$ it is $28$, between $5^2$ and $6^2$.
Thus the quadratic is irreducible for every $a\geq2$.

\begin{remark}
At $a=1,b=2$ the cubic is $(x-3)(x-6)(x-8)$, although the full graph
has irrational antisymmetric eigenvalues $(13\pm\sqrt{33})/2$.
A direct construction of its $10$ vertices gives
\[
\det(xI-L)=x(x-10)(x-9)^3(x-7)^2(x-4)(x^2-13x+34).
\]
Thus a noninteger cubic root is a sufficient condition for graph
nonintegrality; the converse cannot be used without the other blocks.
\end{remark}

\section{Remaining questions}

For repeated exponents $a\geq10$, the next question is whether the
symmetric cubic $f_{a,b}$ always has a noninteger root on the admissible
square-discriminant pairs with
\[
1\leq b\leq\left\lfloor\frac{2(a-4)(a-2)}{3a+4}\right\rfloor.
\]
Their smaller factor satisfies $d\geq3a+4$, with index $j\geq3$.
The two divisor criteria give finite problems for each fixed $a$ or each
fixed $j$, but the resulting family over arbitrary parameters is not settled.
The cases with three pairwise distinct exponents, and nonsquarefree
exponent vectors with more than three prime factors, require further
arguments. The results here therefore do not establish nonintegrality
for every integer with at least three distinct prime factors.

A useful direction is to combine the arithmetic conditions on $(d,e)$
with a sign or congruence obstruction for the symmetric cubic. The
distinction between the two blocks is essential: an integral
antisymmetric block occurs on the boundary family, while an integral
symmetric cubic occurs at $(1,1,2)$.

\begin{thebibliography}{9}
\bibitem{nikandish}
R. Nikandish,
\emph{Laplacian Spectrum of the Ideal Intersection Graph of $\mathbb Z_n$:
Generalized Joins and a Schur--Weyl Reduction}, Zenodo (2026).
\href{https://doi.org/10.5281/zenodo.23134979}{doi:10.5281/zenodo.23134979}.
\bibitem{intersection}
I. Chakrabarty, S. Ghosh, T. K. Mukherjee and M. K. Sen,
\emph{Intersection graphs of ideals of rings},
Discrete Mathematics \textbf{309} (2009), 5381--5392.
\href{https://doi.org/10.1016/j.disc.2008.11.034}{doi:10.1016/j.disc.2008.11.034}.
\end{thebibliography}
\appendix
\section{Exact verification and reproducibility}
\label{sec:verification}

The infinite conclusions use the written proofs and the finite certificates
to which those proofs reduce the problem. Computations use integer or
rational arithmetic; no floating-point spectrum or Lean formalization is
claimed for these Laplacian results. SymPy~1.14.0 and Python~3.10 or later
are used, with source hashes in the saved results.

The \href{https://github.com/the-omega-institute/ideal-intersection-laplacian}{project repository}
contains every checker, exact certificate and reproduction command~\cite{repository}.
The build guide gives the commands by filename. The complete per-checker
catalogue is retained in \path{paper/sections/verification-details.tex};
the optional detailed build includes it without changing the proofs or
saved certificates. The principal verification scopes are as follows.

\subsection{Graph embeddings and repeated exponents}

Direct ideal adjacency validates $LU=UB$, independently of the support
construction, on $14$ squarefree and $(1,1,k)$ cases with $44\,294$
vertex pairs. The full $34$-vertex polynomial at $(2,2,3)$ is checked.
Three boundary graphs with $34,174,548$ vertices cover $165\,490$ pairs
and verify both repeated-exponent embeddings and the quotient lifting.
The $548$-vertex full characteristic polynomial is not computed.

The $90$ cubic diagnostics have exact rational intervals and independent
Sturm/sign checks. The all-$b$ proofs use the exhaustive divisor reduction,
not extrapolation from those diagnostics. All divisors of $24,108,320$
for smaller-factor indices $1,2,3$ are exhausted; eight positive-$b$
cubics have modular certificates checked by separate integer Horner
arithmetic. Higher indices use the uniform sign proof.
The separate ten-vertex graph for $(1,1,2)$ confirms why the symmetric
cubic alone is not an integrality characterization.

\subsection{Small minima and complete finite domains}

The general six-support quintic agrees with the reflected source quotient
and weighted symmetry. Thirty-five requested fully distinct triples and
direct ideal graphs with $22,58$ vertices check the construction. Endpoint
identities and positive shifted coefficients support the written unit,
minimum-two and uniform-cutoff proofs.

At minimum three, the proved cutoff leaves precisely $325$ pairs
$4\leq b<c<30$. All integer quintics, endpoint signs and intervals are
checked; the first endpoint has $257$ positive and $68$ negative values,
with no zero. This complete certificate is part of the minimum-three proof.
At minima four through seven, all $27\,562$ ordered triples in the derived
region are exhausted. A separate fraction-free Bareiss implementation
reconstructs every discriminant by a $9\times9$ integer Sylvester
determinant, checking exact divisions, consecutive-square brackets and
domain completeness. Together with the written cutoff, this proves
Corollary~\ref{cor:minimum-seven}; the minimum itself is not extrapolated.

\subsection{Inertia and endpoint surfaces}

Symbolic block, Schur-complement, determinant and positivity identities
support the uniform low root and balanced region. Six secular/Sturm
fixtures cover the possible nonsingular sign combinations; nine direct
quotient/Sturm fixtures include equal exponents and zero-diagonal boundaries.
The four bounded-gap families also have complete $38$-residue modular
certificates. These finite fixtures check formulas; the infinite conclusions
use the written comparison and inertia arguments.

For the middle-exponent bound, all $132$ cubics over $66$ pairs
$4\leq a<b\leq15$ are independently reconstructed. Integer trial division
exhausts $8\,527$ positive divisors, and separate Horner evaluation checks
all $7\,585$ above $b$. No upper bound on $c$ is imposed. This is the
complete certificate underlying Theorem~\ref{thm:endpoint-reduction}.
The linear cutoff uses the general endpoint cubic and five positive
coefficient expressions. Both known repeated-exponent endpoint zeros are
positive controls. The CRT scope note uses six substitution identities,
two positive boundary shifts and two integer seed determinants; its
all-moduli claim is the written polynomial-congruence argument.

\subsection{Arithmetic and shifted-root cases}

Scaling, all parity patterns and common residues for primes $3,5,7$
are checked coefficientwise, with $70$ quadratic-residue evaluations and
separate integer Horner arithmetic. The all-even scaling contradiction
uses the written valuation proof. Complete endpoint tables modulo two
and three include zero residues and zero polynomials; $70$ independent
integer determinant evaluations check the endpoint table.

For the cases in Appendix~\ref{sec:two-adic-lifting}, direct and reflected
quintics and all exponent permutations agree. Every shifted divisibility
and binary-cubic identity is checked coefficientwise for arbitrary shifts.
Chosen full quintics and derivatives are independently reconstructed by
integer determinants and principal cofactors, including compatible controls
and both known endpoint-zero examples. The finite residue and parity tables
support the identities; the infinite conclusions use the written root
factorizations, rather than a scan of exponents.

Reza's higher-modulus script was run unchanged; the independent verifier
checks all $28\,032$ residue triples with $56\,064$ Bareiss determinants.
The coefficientwise period-eight proof explains why the fixed divisor-$32$
tables add no classes at larger exponent moduli. His derivative script's
$384$ representatives are checked by $16\,380$ principal-minor determinants,
including the additional controls. Zero has infinite valuation, and full
valuations need not be constant on an exponent residue class. The complete
catalogue retains these diagnostic limits, coefficient counts, fixture
counts and exact source/certificate mappings.

\subsection{Endpoint-one spectral location}

Theorem~\ref{thm:endpoint-one-spectrum} and
Corollary~\ref{cor:endpoint-one-divisor-window} use written proofs without
a finite base. The
\href{https://github.com/the-omega-institute/ideal-intersection-laplacian/blob/d09cfdcc08f2a0b6c5b3abbd7d7a30bacc9030b4/scripts/check_endpoint_one_minimum_root.py}{minimum-root checker}
reconstructs the generic endpoint polynomial, comparison and equality
factorization, all twelve positive terms and the determinant identities.
The
\href{https://github.com/the-omega-institute/ideal-intersection-laplacian/blob/fd02ed726d3ada4bbd5e401e8818739f3bddc371/scripts/check_endpoint_one_pair_sum_window.py}{pair-sum checker}
verifies the four-support block, characteristic factor, discriminant and
strict sign identity. The established controls $(8,8,105)$, $(9,9,136)$
and $(20,20,741)$ leave $5,5,13$ divisors in the smallest-root windows;
all $23$ integer Horner evaluations are nonzero. Rational factorization
agrees, and exact Sturm counts give one quartic root in each window and
none on $[0,a]$. These controls already follow from the repeated-exponent
theorem. The source-hashed certificates reproduce byte-for-byte; their
finite diagnostics supplement the written bounds.

Proposition~\ref{prop:endpoint-one-equality-spectrum} also has a written
proof without a finite base. The
\href{https://github.com/the-omega-institute/ideal-intersection-laplacian/blob/d3fa3478193f34777a88b970da8f0763ce177c38/scripts/check_endpoint_one_equality_root.py}{equality-index checker}
verifies the principal factors, weighted symmetry and nonzero coupling.
The
\href{https://github.com/the-omega-institute/ideal-intersection-laplacian/blob/9e98266eb279a31b7282cea17fca3f5af6d9410f/scripts/check_endpoint_one_equality_second_root.py}{equality pair-sum checker}
verifies both positive identities at $a+b$. The three fixed controls
$(4,4,20)$, $(9,12,159)$ and $(9,30,951)$ are all off endpoint one.
After removing root $b$, exact Sturm counts place one positive root below
$a+b$ and three above it. These are diagnostics of the unconditional
statement, not integer endpoint solutions or a finite base.

\ifdefined\DetailedVerification
\subsection{Complete checker catalogue}
\begin{enumerate}
\item \path{scripts/check_integrality_obstructions.py} constructs ideals
as exponent tuples and tests adjacency by coordinatewise maxima, independently
of the support description. It verifies $LU=UB$ for the displayed
invariant embeddings, for squarefree $t=3,\ldots,8$ and $(1,1,k)$ with
$k=1,\ldots,8$: $14$ cases and $44\,294$ vertex pairs. It also checks
the square-gap inequalities and principal-minor witnesses.

\item \path{scripts/check_remaining_case.py} constructs all $34$ vertices
for $(2,2,3)$ and verifies its full characteristic polynomial and the
class decomposition in exact arithmetic. The shifted symmetric cubic is
$x^3-80x^2+2099x-18052$; its reduction modulo $5$ has no root.

\item \path{scripts/check_repeated_exponent_family.py} checks the general
quotient polynomial, both boundary sign identities, the positive-coefficient
expansion, and the integral antisymmetric eigenvalues symbolically.
Direct graphs for $(2,2,3),(3,3,10),(4,4,21)$ have $34,174,548$ vertices,
covering $165\,490$ vertex pairs. Exact neighbor counts verify both blocks
and their lifting. The check computes the seven-support quotient polynomial;
it does not compute the full $548\times548$ characteristic polynomial.

\item \path{scripts/check_aa_b.py} factors the $90$ diagnostic cubics
over $\mathbb Q$ and uses exact rational root isolation for a root between
consecutive integers. Independent Sturm counts and endpoint signs were
checked for every saved interval. Its modular search is restricted to
$2,3,5,7,11,13,17,19,23,29,31$; absence of a listed witness is not a proof
that no prime works for that cubic.

\item \path{scripts/check_fixed_repeated_exponents.py} reconstructs
the cubic from $B$, checks the discriminant, product, cutoff and reducible
factor identities, and enumerates every factor pair for $a=2,3,4$.
It verifies the admissibility conditions and all modular witnesses in
Tables~\ref{tab:divisors} and~\ref{tab:cubics}. The general congruence
and factor inequality are proved in Section~\ref{sec:bounds}.
\item \path{scripts/check_second_cutoff.py} checks the polynomial
remainder modulo $a+2$, enumerates every divisor of $24$, and verifies
the four root-free cubic reductions settling all $d=a+2$ cases.
It also checks the second-cutoff identities and the two nonsquare
discriminants completing $a=5,6$. This finite certificate follows the
proved divisor reduction; it is not a parameter-range search.
\item \path{scripts/check_factor_index_reduction.py} verifies the general
fixed-index divisibility and reconstruction identities, exhausts the
divisors of $K_2=108$, and certifies the unique positive-$b$ index-$2$
cubic modulo $13$. It checks the third-cutoff identities and the four
remaining discriminants for $a=7,8,9$. These finite checks exhaust the
proved candidate sets, with no arbitrary parameter cutoff.
\item \path{scripts/check_repeated_exponent_completion.py} reconstructs
the general cubic from the $4\times4$ quotient, checks both sign-interval
identities, and verifies the fourth-cutoff and strictly positive gap
identities. It exhausts every divisor of $24,108,320$ for indices
$1,2,3$, reproducing all eight positive-$b$ modular cubic certificates
with separate integer Horner arithmetic. The infinite remaining range is
settled by the sign proof of Theorem~\ref{thm:repeated}, not by a scan.
\item \path{scripts/check_six_support_quotient.py} verifies the general
six-support and complement quotients, their weighted symmetry and reflected
quintics, and the two determinant values used in Theorem~\ref{thm:unit}.
Exact rational factorization covers all thirty-five requested triples
$1\leq a<b<c\leq7$, with every matrix and quintic cross-checked against
\path{scripts/check_fully_distinct.py}. All saved rational intervals
pass independent Sturm counts and endpoint-sign checks.
Direct ideal adjacency verifies the quotient for
$(1,2,3)$ and $(2,3,4)$, using $22$ and $58$ vertices and $312$
representative-to-vertex pairs. The infinite unit-exponent conclusion
uses the written endpoint proof.
\item \path{scripts/check_minimum_two.py} verifies the complement
endpoint identities at $a=2$, the symmetric expression in $b+c$ and $bc$,
and both positive-coefficient expansions that exhaust the infinite range.
It reconstructs the single exceptional quintic at $(2,3,4)$ and checks
its rational endpoint values with separate Horner evaluation and a Sturm
count. No new parameter range is enumerated.
\item \path{scripts/check_distinct_tail.py} reconstructs the general
six-support complement quintic and checks its cubic first endpoint,
the uniform cutoff identity and positive remainder, and both
minimum-three second-endpoint expansions. The cutoff itself derives
the finite domain $4\leq b<c<30$, containing exactly $325$ pairs.
For every pair the checker reconstructs the integer quotient quintic,
matches its symbolic specialization, and checks the displayed root
interval using separate rational Horner arithmetic. Separate cubic
Horner evaluation confirms all first-endpoint values: $257$ positive,
$68$ negative and no zero. The $(3,4,5)$ rational exception also has a
Sturm count of one in $(1,3/2)$. The minimum-three theorem uses this
complete finite certificate; the general cutoff uses the written
positive-coefficient argument.
\item \path{scripts/check_low_spectrum.py} reconstructs the weighted
symmetric complement quotient and checks its block form, Schur complement,
determinant identity and ordered-parameter derivative. For the four
bounded-gap families it verifies the upper-inertia expressions and both
endpoint factorizations, positive shifted first-endpoint coefficients,
and all $38$ root-free modular residues with separate integer Horner
arithmetic. Eight chosen representatives, the four families at $a=4,20$,
have independently reconstructed integer quotient quintics and Sturm
counts of two roots in $(0,3)$; at $a=20$ both lie in $(2,3)$.
These examples check the formulas. The infinite families use the written
inertia proof and the complete modular exclusions, with no parameter scan.
\item \path{scripts/check_open_region2.py} runs the coauthor-requested
domain $4\leq a\leq7$, $a<b<c<4a^2-2a$. It constructs the generic
$H$ quotient determinant once and specializes its integer coefficients
for every one of the $27\,562$ triples. Every discriminant is nonsquare.
Original-quotient and complement endpoint signs are reported separately;
the general reflection identity is checked symbolically. Sixteen
representatives, including every first complement sign-pattern witness,
have separately reconstructed integer matrix characteristic polynomials,
complement endpoints and exact rational factorizations. The compatibility
entrypoint \path{scripts/check_open_region.py} runs the same corrected code.
\item \path{scripts/check_open_region_certificate.py} independently
reconstructs every one of the $27\,562$ discriminants as a $9\times9$
integer Sylvester determinant, using a separate fraction-free Bareiss
implementation. It checks every exact division, every strict consecutive-square
bracket and the complete ordered domain. This is the exhaustive finite
certificate used in Corollary~\ref{cor:minimum-seven}, combined with the
written cutoff, rather than extrapolation to arbitrary minimum exponent.
\item \path{scripts/check_mixed_inertia.py} verifies the comparison
determinant, negative Rayleigh value, equal-exponent case, positive
comparison gap and balanced-span identities symbolically in a rational
$3\times3$ congruence. Six exact secular/Sturm fixtures cover every
nonsingular diagonal-sign/secular-sign combination possible here.
Nine separately reconstructed quotient/Sturm fixtures include equal
exponents, two zero-diagonal boundaries, a strongly unbalanced triple
and chosen spans up to $1000$. The finite fixtures validate formulas;
Theorem~\ref{thm:uniform-low-root} and
Corollaries~\ref{cor:balanced-interval}--\ref{cor:balanced-span}
use the written comparison and positivity proofs, with no parameter scan.
\item \path{scripts/check_arithmetic_obstructions.py} verifies symbolic
integer quotient scaling, the equal-exponent quintic factorization and
its quadratic discriminant. It checks all eight parity patterns and
every nonzero common residue for primes $3,5,7$, with all $70$ quadratic
residue evaluations, including the binary factor, cross-checked by
integer Horner arithmetic. Six direct quotient fixtures verify scaling,
low-root and modular witnesses, and the existing boundary triple
$(9,9,136)$ checks that an integer endpoint is compatible with other
noninteger roots. The infinite conclusions use the written divisibility
and finite-field proofs, with no exponent scan.
\item \path{scripts/check_endpoint_reduction.py} reconstructs the general
endpoint-two cubic in the middle exponent and the two endpoint constant
terms in the largest exponent. It verifies both rational-cutoff identities
and five positive shifted-coefficient expressions. Three chosen even
gcd-two fixtures below the earlier quadratic cutoff have direct integer
quotient, independent Horner and exact Sturm checks. The existing repeated
triples $(9,9,136)$ and $(10,10,12)$ check the endpoint divisibility
conditions and retain nonlinear rational factors; the latter checks the
strict linear bound. The infinite bound and even-tail corollary use written
positivity and divisibility proofs, without an exponent scan or divisor-set
enumeration. The companion \path{scripts/check_endpoint_congruence_scope.py}
checks six substitution identities, two positive boundary shifts and the two
seed endpoint zeros by integer six-by-six determinants. Its all-moduli
conclusion uses the written polynomial-congruence argument; it does not
assert actual zeros or integer spectra for the constructed candidates.
\item \path{scripts/check_even_exponent_congruence.py} verifies integer
quotient scaling, the independent binary block characteristic polynomial,
and the coefficientwise congruences modulo two and thirty-two for the
scaled quintic. Three chosen gcd-two triples inside the preceding even
linear bound have separately reconstructed integer quotients, rational
factorizations and independent integer Horner values at one and three.
The infinite congruence-family theorem follows from the written divisibility
contradiction, without an exponent scan.
\item \path{scripts/check_endpoint_surfaces.py} is Reza's requested
diagnostic, corrected to use $|V(H)|$ in the complement reflection and
to include $b=15$. Exact rational linear-factor extraction finds no
endpoint roots for $4\leq a<b\leq15$, $c>b$.
\path{scripts/verify_endpoint_surfaces.py} independently reconstructs
the direct disjoint-support complement matrix and all $132$ endpoint
cubics over the $66$ pairs. It exhausts $8\,527$ positive constant-term
divisors by integer trial division and checks all $7\,585$ divisors
above $b$ with separate integer Horner arithmetic. Two previously
settled repeated-exponent positive controls recover their endpoint roots
and match full quotient factorizations. This complete finite certificate,
combined with the uniform low-root theorem and earlier cases, proves
Proposition~\ref{cor:middle-fifteen}. There is no upper cutoff on $c$ or
extension beyond the requested pair range.
\item \path{scripts/check_endpoint_residues.py} reconstructs the general
endpoint cubics and records all $26$ pair/endpoint rows for primes two
and three, including zero residues and zero polynomials. At all $35$
positive residue representatives, a separate $720$-term integer
determinant expansion checks both endpoint values, for $70$ checks.
The eight parity reductions and five simultaneous modulo-three exclusions
are verified against direct quotient characteristic polynomials.
Two selected fully distinct triples have independent endpoint determinants
and exact Sturm counts; two known repeated-exponent controls retain their
endpoint zeros. The infinite classes in
Proposition~\ref{thm:modulo-three-endpoints} use the written low-root proof.
This is a residue enumeration, not an exponent search.
\item \path{scripts/check_mixed_parity_congruence.py} checks the two
mixed binary reductions and four initial coefficient-divisibility identities.
It verifies all $558$ coefficients in twelve substitutions for the six
rows of Proposition~\ref{thm:mixed-parity-congruence}. Six positive residue
representatives and two selected fully distinct triples have $16$
independent $720$-term integer determinant checks at one and three,
matching separate integer Horner values. The fixtures retain nonlinear
rational factors. Complete lower-modulus tables give compatible linear
root multisets for all $12$ and $80$ unordered mixed-parity exponent
classes modulo four and eight, respectively; this compatibility is not
an integrality claim. The infinite theorem follows from the written
three-odd-root divisibility lemma.
\item \path{scripts/check_odd_root_distribution.py} verifies the generic
quintic's six exponent permutations and the coefficientwise identities
in Proposition~\ref{thm:odd-root-distribution}: $160$ coefficients for the
shifted polynomial's divisibility by eight, $160$ for its reduction
modulo two, $42$ endpoint coefficients modulo thirty-two and $51$
derivative coefficients modulo sixteen. Four selected fixtures have
$28$ independently expanded integer six-by-six determinants, including
six interpolation points and the zero-root check per fixture. A further
$24$ five-by-five principal cofactors independently reconstruct the
derivative at one. The reconstructed quintics and derivatives agree;
the infinite theorem uses the written odd-root distribution argument.
\item Reza's \path{scripts/check_2adic_higher_moduli.py} was run unchanged
at exponent moduli eight, sixteen and thirty-two. The independent
\path{scripts/verify_2adic_higher_moduli.py} checks all $28{,}032$ ordered
mixed-parity residue triples using $56{,}064$ integer Bareiss determinants
for the two endpoints. Twelve coefficientwise identities, with $560$
coefficients, establish coordinate period eight modulo thirty-two.
The complete $384$ base rows therefore encode the higher tables.
The endpoint-divisibility test excludes $36$, $288$ and $2304$ classes,
respectively. These are finite residue tables, and their survivors
do not assert integer endpoint zeros or integral spectra.
\item Reza's \path{scripts/check_higher_derivatives.py} was run unchanged
on its $384$ canonical representatives. Its interpretation was then
corrected: zero has valuation infinity, and full valuations need not be
constant on an exponent residue class. The independent
\path{scripts/verify_higher_derivatives.py} reconstructs all four values
at one by principal integer minors. Two new fixtures at three and four
lift controls at one bring the total to $16{,}380$ minor determinants.
Twelve coefficientwise identities determine the odd-root distribution
modulo four; $36$ further identities establish coordinate period eight
for the derivative residues at one and three modulo $16$, $8$ and $4$.
The second- and third-derivative thresholds add no exclusions with this
root information. A coefficientwise identity for $h_C(3)/16$ proves
Proposition~\ref{thm:odd-root-sum}; its $64$ role-ordered classes modulo
sixteen display $48$ exclusions, including $16$ beyond the earlier
divisor-$32$ test. Full joint representative patterns are saved separately.
\item \path{scripts/check_odd_root_lift.py} matches the direct complement
quintic to the reflected source and checks all six exponent permutations.
All $185$ coefficients of the second shifted polynomial are divisible
by $64$, and $185$ difference coefficients verify its binary cubic
identity for arbitrary integer shifts. Eight shift-parity patterns show
the irreducible quadratic in the excluded cases. Four fixtures, two
new exclusions and two compatible controls, independently reconstruct
their full quintics by $28$ integer determinant expansions, including
the zero-root checks. The $64$ additional role-ordered exclusions
modulo thirty-two follow from the written congruence and its class count;
no exponent or modulus range scan is needed for
Proposition~\ref{thm:odd-root-lift}.
\item \path{scripts/check_clustered_odd_roots.py} matches the direct
and reflected generic quintics, all exponent permutations, and the
exact factorization of $h_C(b)$. It checks $84$ difference coefficients
for the derivative congruence modulo sixty-four, and $249$ scaled
coefficients and $249$ binary-difference coefficients for each of the
two third-shift identities. Four derivative residues and sixteen binary
shift-parity patterns display the resulting necessary conditions.
Eight fixtures independently reconstruct full quintics with $56$ integer
six-by-six determinant expansions; another $48$ five-by-five principal
cofactors reconstruct their derivatives at $b$. Two fixtures have
$h_C(b)=0$ and satisfy the derivative divisor but fail the third-shift
condition. Four compatible controls retain nonlinear rational factors.
The arithmetic conditions in Proposition~\ref{thm:clustered-odd-roots} use
the root-clustering proof in Appendix~\ref{sec:two-adic-lifting}, with no
exponent or modulus range scan.
\item \path{scripts/check_mixed_endpoint_roots.py} matches the direct
and reflected quintics and all exponent permutations. For the three
mixed-parity patterns, coefficientwise checks verify the binary quintic,
divisibility of the shifted polynomial by eight, and its binary reduction
to $(y+1)^3$. A further $51$ coefficients give $h_C(2)\equiv4\pmod8$
for residue roles $(3,0,0)$. Eight full quintics are independently
reconstructed by $56$ integer six-by-six determinant expansions;
exact Sturm counts check the low-root fixtures. Two previous clustered-root
controls pass all prior necessary conditions but are excluded by the
linear tail. Both known endpoint-zero controls are retained.
Proposition~\ref{thm:one-odd-endpoints} and
Proposition~\ref{cor:mixed-linear-tail} follow from the written root-distribution
argument and the preceding low-root and endpoint-two proofs.
No exponent or modulus range is scanned.
\item The pinned PR9 \path{check_endpoint_two_middle_tail.py} verifies
the full endpoint Schur determinant, the diagonal factorization and
the middle-bound positive expansion, with five selected quotient/Horner/Sturm
fixtures and the repeated endpoint-zero control. Its unbounded statement
uses the written positive identity.
\item The pinned PR9 \path{check_endpoint_two_sharp_maximum.py} verifies
the inverse interval substitution, both nonnegative expressions and every
coefficient of the $84$-term residual in
Lemma~\ref{lem:endpoint-two-rectangle}. All $16$ residual vectors agree
with Appendix~\ref{sec:endpoint-two-identities}; eight direct fixtures
and two below-cutoff controls supplement the written proof.
\item At PR9, \path{check_endpoint_two_completion.py} checks the
$44$- and $112$-term identities and all $32$ coefficient vectors.
It exhausts exactly $8\,658$ necessary fully distinct triples, with
independent integer-Horner and $6\times6$ Bareiss evaluations in every
case. The complete CSV and JSON reproduce byte-for-byte. Seven direct
quotient/Horner/Sturm fixtures and the repeated control supplement the
proof; earlier minimum-through-seven certificates remain dependencies.
\item At PR9, \path{check_endpoint_one_equality_cubic.py} verifies
the explicit cubic's generic quotient identity and discriminant by an
independent $5\times5$ Sylvester determinant. Its $61$-term expansion,
divisor/quadratic identity and abstract square-discriminant counterexample
are exact checks, with no exponent scan or point enumeration.
\end{enumerate}

The symbolic and characteristic-polynomial checks use SymPy~1.14.0;
the first checker uses the Python standard library.
Saved results contain source hashes. The $10$-vertex polynomial in the
preceding remark was independently checked by a direct ideal-graph
construction; it is also reproducible by the command in the manuscript
build guide. No floating-point eigenvalue calculation is used as
integrality evidence.

\fi

\section{Mixed-parity case analyses and two-adic lifting}
\label{sec:two-adic-lifting}

Theorem~\ref{thm:mixed-residue-classes} collects the following cases.
All residue-role permutations are included. Each row records either a
nonintegral class or a necessary condition for an integer spectrum.
\begin{center}
\begin{tabular}{p{3.2cm}p{1.6cm}p{8.8cm}}
\toprule
Residue roles & Modulus & Conclusion or necessary condition\\
\midrule
Six classes listed in the main theorem & $8$ & Nonintegral by the divisor-$32$ endpoint table below.\\
$(3,3,2)$ & $4$ & Nonintegral: endpoint and derivative requirements are incompatible.\\
$(3,0,0)$ & $4$ & Nonintegral: one is excluded by root distribution and $h_C(2)\equiv4\pmod8$.\\
$(1,3,2)$ & $4$ & $a+c\equiv b(b+2)\pmod{128}$, $b\equiv11,15\pmod{16}$ and the stated binary-cubic parity condition.\\
$(1,3,2)$, $(3,3,0)$ & $4$ & For minimum $m\geq4$, $h_C(2)=0$ and maximum $M<R(m)$.\\
\bottomrule
\end{tabular}
\end{center}

The complete identities, intermediate conditions, examples and diagnostic
boundaries are retained below. The variables identify residue roles rather
than size order; no additional refinement or search is introduced.

\subsection{Endpoint divisibility and odd-root distribution}

\begin{lemma}
\label{lem:three-odd-roots}
Let $f\in\mathbb Z[x]$ be monic of degree five with
$f(x)\equiv x^2(x+1)^3\pmod2$. If all its roots are integers,
then $32\mid f(1)$ or $32\mid f(3)$.
\end{lemma}

\begin{proof}
The binary factorization forces precisely three odd roots and two even
roots, counted with multiplicity. At least two odd roots share a residue
$r\in\{1,3\}$ modulo four. In $f(r)$, their two factors are divisible
by four and the third odd-root factor is divisible by two. The even-root
factors are odd. Therefore $32\mid f(r)$, including when $f(r)=0$.
\end{proof}

\begin{proposition}
\label{thm:mixed-parity-congruence}
Let $a,b$ be the two odd exponents and $c$ the even exponent of a positive
three-prime exponent triple. Each row in the following table, including
exchange of $a,b$, implies that its ideal intersection graph is nonintegral.
\begin{center}
\begin{tabular}{ccc|cc}
\toprule
$a\bmod8$ & $b\bmod8$ & $c\bmod8$ & $h_C(1)\bmod32$ & $h_C(3)\bmod32$\\
\midrule
1&3&2&8&16\\
1&7&2&24&16\\
3&7&2&16&8\\
3&5&6&8&16\\
3&7&6&16&24\\
5&7&6&24&16\\
\bottomrule
\end{tabular}
\end{center}
\end{proposition}

\begin{proof}
For either mixed exponent parity,
$h_C(x)\equiv x^2(x+1)^3\pmod2$. Substituting
$(a,b,c)=(8u+r_1,8v+r_2,8w+r_3)$ in the general quotient quintic gives
the last two columns: every nonconstant coefficient in $u,v,w$ is
divisible by $32$. Both entries are nonzero modulo $32$, contradicting
Lemma~\ref{lem:three-odd-roots} if all quotient roots were integers.
The resulting noninteger real complement eigenvalue has a noninteger
graph lift.
\end{proof}

For example, $(10,17,19)$ and $(14,19,21)$ meet these criteria with
gcd one and unequal $2$-adic valuations. There is no minimum or ratio
bound in the argument. The weaker divisibilities $4\mid h_C(2t)$
and $8\mid h_C(2t+1)$ hold identically for both mixed exponent parities;
they alone do not supply the contradiction. The distribution of the
three odd roots modulo four is the additional constraint.

\begin{lemma}
\label{lem:odd-root-distribution}
Let $f\in\mathbb Z[x]$ be monic of degree five with
$f(x)\equiv x^2(x+1)^3\pmod2$. Suppose
$F(y)=f(2y+1)/8$ belongs to $\mathbb Z[y]$ and
$F(y)\equiv y^3\pmod2$. If all roots of $f$ are integers,
then its three odd roots are congruent to one modulo four,
$64\mid f(1)$ and $16\mid f'(1)$.
\end{lemma}

\begin{proof}
Write the two even roots as $\eta_1,\eta_2$ and the three odd
roots as $2r_1+1,2r_2+1,2r_3+1$, with multiplicities. Then
\[
 F(y)=\prod_{j=1}^2(2y+1-\eta_j)\prod_{i=1}^3(y-r_i).
\]
The first two factors reduce to one modulo two. Hence
$\prod_i(y-r_i)\equiv y^3\pmod2$, which forces every $r_i$ to be even.
Each odd-root factor in $f(1)$ is therefore divisible by four,
giving $64\mid f(1)$. Every term in the product-rule expansion of
$f'(1)$ retains at least two such factors, giving $16\mid f'(1)$.
The argument includes coincident roots and zero factors.
\end{proof}

\begin{proposition}
\label{thm:odd-root-distribution}
Every positive three-prime exponent triple that is a permutation of
$(3,3,2)$ modulo four has a nonintegral ideal intersection graph.
\end{proposition}

\begin{proof}
Relabel the odd exponents as $a=4u+3$, $b=4v+3$ and the even
exponent as $c=4w+2$, where $u,v,w\geq0$. Expansion of the
general complement quintic gives the coefficientwise identities
\begin{align*}
 h_C(2y+1)/8&\in\mathbb Z[u,v,w,y], &
 h_C(2y+1)/8&\equiv y^3\pmod2,\\
 h_C(1)&\equiv16(u+v)\pmod{32}, &
 h_C'(1)&\equiv8(u+v+1)\pmod{16}.
\end{align*}
Its binary reduction is $x^2(x+1)^3$. If all five roots were integers,
Lemma~\ref{lem:odd-root-distribution} would give $64\mid h_C(1)$,
so the endpoint congruence would require $u+v$ even. The derivative
divisibility $16\mid h_C'(1)$ would require $u+v$ odd, a contradiction.
Thus there is a noninteger real complement eigenvalue, whose graph
lift is noninteger.
\end{proof}

This theorem contains the two rows $(3,7,2)$ and $(3,7,6)$ of
Proposition~\ref{thm:mixed-parity-congruence} and also covers odd exponents
sharing the same residue modulo eight. For example, $(10,19,27)$ and
$(14,23,39)$ have odd residues $(3,3)$ and $(7,7)$ modulo eight,
respectively. Neither theorem requires a low-root interval or an
exponent-size bound.

For the fixed necessary divisor $32$, increasing the exponent modulus
alone gives no further exclusions. On all mixed-parity triples, the
values of both $h_C(1)$ and $h_C(3)$ modulo $32$ already have period
eight in each exponent. This follows coefficientwise by substituting
the two parity representatives $(2u+1,2v,2w)$ and $(2u+1,2v+1,2w)$
and shifting one variable by $4z$; symmetry covers their permutations.
Thus the higher-modulus tables are lifts of the six classes in
Proposition~\ref{thm:mixed-parity-congruence}. The additional exclusions in
Proposition~\ref{thm:odd-root-distribution} use the distribution of the roots
and the derivative, beyond those two endpoint divisibilities.

\begin{proposition}
\label{thm:clustered-odd-roots}
Let $a,b,c$ be positive exponent residue roles with
$a\equiv1$, $b\equiv3$ and $c\equiv2\pmod4$.
If their ideal intersection graph is Laplacian integral, then
\[
 a+c\equiv b(b+2)\pmod{128},\qquad b\equiv11\text{ or }15\pmod{16}.
\]
Moreover, put $u=(a-1)/4$ and $t=(a+c-b(b+2))/128$. Then
\[
 \begin{cases}
 t+u+v\equiv1\pmod2,& b=16v+11,\\
 t+v\equiv1\pmod2,& b=16v+15.
 \end{cases}
\]
If $m=\min\{a,b,c\}\geq4$ and $M=\max\{a,b,c\}$, then also
\[
 h_C(2)=0,\qquad M<7m-16+\frac{40}{m+2}.
\]
Every violation under the stated hypotheses proves nonintegrality;
all exponent permutations are included.
\end{proposition}

\begin{proof}
Proposition~\ref{thm:odd-root-lift} in Appendix~\ref{sec:two-adic-lifting}
forces all three hypothetical integer odd quotient roots to be $b$ modulo
eight. Thus $512\mid h_C(b)$ and $64\mid h_C'(b)$.
The exact identity
\[
 h_C(b)=a^2b^2c^2\bigl(a+c-b(b+2)\bigr)
\]
has prefactor of $2$-adic valuation two, giving the sum condition modulo
$128$, including when $h_C(b)=0$. The derivative congruence and the binary
cubic calculation in Appendix~\ref{sec:two-adic-lifting} give the restriction
on $b$ and the two parity conditions.

Since $b\equiv3\pmod4$, one cannot be a quotient root under the
integer-spectrum assumption. For $m\geq4$, the positive root in $(0,3)$
from Theorem~\ref{thm:uniform-low-root} must therefore equal two.
Proposition~\ref{thm:endpoint-two-bound} gives the strict linear upper bound.
\end{proof}

The arithmetic conditions imply the earlier sum congruences
$a+c\equiv15\pmod{16}$ and $a+c+4b\equiv27\pmod{32}$.
Appendix~\ref{sec:two-adic-lifting} gives their staged derivation and the
supporting derivative thresholds. The theorem collects these arithmetic
conditions together with the endpoint-two restriction.
For example, $(29,27,754)$ and $(45,47,2258)$ have $h_C(b)=0$ and
$64\mid h_C'(b)$ but fail the final parity test. The compatible triples
$(25,27,758)$ and $(29,27,882)$ instead fail the linear upper bound.
Satisfying the arithmetic conditions alone does not establish integrality.

\begin{lemma}
\label{lem:mixed-endpoint-one}
For positive exponent triples that are permutations of $(3,0,0)$,
$(1,3,2)$ or $(3,3,0)$ modulo four, an integer complement quotient
spectrum would force every odd root to be three modulo four.
Consequently one cannot be a quotient root under this assumption.
\end{lemma}

\begin{proof}
For each displayed residue pattern, write
$(a,b,c)=(4u+r_a,4v+r_b,4w+r_c)$. Coefficientwise expansion gives
\begin{align*}
 h_C(x)&\equiv x^2(x+1)^3\pmod2,\\
 F(y)=h_C(2y+1)/8&\in\mathbb Z[u,v,w,y], &
 F(y)&\equiv(y+1)^3\pmod2.
\end{align*}
If every root were an integer, there would be three odd roots $2r_i+1$
and two even roots $\eta_j$, counted with multiplicity. In
\[
 F(y)=\prod_{j=1}^2(2y+1-\eta_j)\prod_{i=1}^3(y-r_i),
\]
the first two factors reduce to one modulo two. Thus every $r_i$ is
odd, giving the asserted residue and exclusion of one. Symmetry of the
quintic includes every permutation of the exponents.
\end{proof}

\begin{proposition}
\label{thm:one-odd-endpoints}
Every positive three-prime exponent triple that is a permutation of
$(3,0,0)$ modulo four gives a nonintegral ideal intersection graph,
without a size or ratio bound.
\end{proposition}

\begin{proof}
For residue roles $(a,b,c)=(4u+3,4v,4w)$, coefficientwise expansion
gives $h_C(2)\equiv4\pmod8$, so two is not a root.
If the minimum exponent is at least four, Theorem~\ref{thm:uniform-low-root}
provides a positive quotient root in $(0,3)$. Under the integer-spectrum
assumption it must be one or two. Lemma~\ref{lem:mixed-endpoint-one}
excludes one and the endpoint congruence excludes two, a contradiction.
The only possible smaller minimum is three, already covered by
Theorem~\ref{thm:minimum-three}. A noninteger complement root has a
noninteger graph lift.
\end{proof}

The triples $(19,36,100)$ and $(35,52,120)$ illustrate this one-odd-exponent
class beyond the minimum-through-seven and middle-through-fifteen
certificates. The result uses the uniform low-root theorem together with
the shifted root distribution, without enlarging either finite certificate.

\subsection{Derivative divisibility}

\begin{lemma}
\label{lem:derivative-thresholds}
Let $f\in\mathbb Z[x]$ be monic of degree five with two even and three
odd integer roots. If exactly $k$ of the odd roots are one modulo four,
then
\[
 2^{3+k}\mid f(1),\qquad
 2^{\max(2,k+1)}\mid f'(1),\qquad
 2^{\max(2,k)}\mid f''(1),\qquad 2\mid f'''(1).
\]
At three, the same statements hold with $k$ replaced by $3-k$.
\end{lemma}

\begin{proof}
At one, the odd-root factors supply $k$ factors divisible by four and
$3-k$ divisible by two; the even-root factors are odd. Each term of
the $j$th derivative is $j!$ times a product retaining $5-j$ root factors.
Removing the odd-root factors with the strongest guaranteed divisibility
gives the stated bounds, including the factor two from $2!$ and $3!$.
At three, the roles of the two odd residue classes are exchanged.
Repeated roots and zero evaluations obey the same divisibilities.
\end{proof}

\subsection{Odd roots modulo four and eight}

\begin{proposition}
\label{thm:odd-root-sum}
Let $a,b,c$ be positive exponents with
$a\equiv1$, $b\equiv3$ and $c\equiv2\pmod4$.
If $a+c\not\equiv15\pmod{16}$, their ideal intersection graph is
nonintegral. All permutations are included, without a size or ratio bound.
\end{proposition}

\begin{proof}
Write $(a,b,c)=(4u+1,4v+3,4w+2)$. Coefficientwise expansion gives
\[
 h_C(2y+1)/8\in\mathbb Z[u,v,w,y],\qquad
 h_C(2y+1)/8\equiv(y+1)^3\pmod2.
\]
As in Lemma~\ref{lem:odd-root-distribution}, integer quotient roots
would force all three odd roots to be three modulo four. Consequently
$64\mid h_C(3)$. A further coefficientwise identity is
\[
 h_C(3)/16\in\mathbb Z[u,v,w],\qquad
 h_C(3)/16\equiv(2v+1)(u+w+1)\pmod4.
\]
Since $2v+1$ is invertible modulo four, this forces
$u+w+1\equiv0\pmod4$, equivalently $a+c\equiv15\pmod{16}$.
Failure of that condition gives a noninteger complement root and graph lift.
\end{proof}

For example, $(14,25,27)$ and $(10,29,31)$ pass the earlier divisor-$32$
test but have $h_C(3)\equiv32\pmod{64}$, contradicting the new requirement.
The sum class $a+c\equiv15\pmod{16}$ remains compatible with this test.
Complete second- and third-derivative checks at one and three give no
additional exclusions with the root information modulo four used here.
In particular $h_C'''(1)\equiv2\pmod4$ for every mixed exponent parity.
Full valuations at a chosen exponent representative are not generally
constant throughout its residue class: $h_C(3)=-2592$ at $(1,3,6)$ but
$h_C(3)=0$ at $(9,3,6)$, although the exponent triples agree modulo eight.

\begin{proposition}
\label{thm:odd-root-lift}
Let $a,b,c$ be positive exponents with
$a\equiv1$, $b\equiv3$ and $c\equiv2\pmod4$.
If $a+c+4b\not\equiv27\pmod{32}$, their ideal intersection graph is
nonintegral. All permutations are included, without a size or ratio bound.
\end{proposition}

\begin{proof}
Suppose all five complement quotient roots are integers.
Proposition~\ref{thm:odd-root-sum} requires $a+c\equiv15\pmod{16}$.
Write
\[
 (a,b,c)=(4u+1,4v+3,16t-4u-2).
\]
Coefficientwise expansion gives
\[
 H(z)=h_C(4z+3)/64\in\mathbb Z[u,v,t,z],\qquad
 H(z)\equiv z^3+vz^2+vz+t+1\pmod2.
\]
The three odd roots are already known to be three modulo four.
Writing them as $4r_i+3$ and the two even roots as $\eta_j$ gives
\[
 H(z)=\prod_{j=1}^2(4z+3-\eta_j)\prod_{i=1}^3(z-r_i).
\]
The first two factors reduce to one modulo two, so the displayed cubic
must split into linear factors over $\mathbb F_2$.
For even $v$, it splits precisely when $t$ is odd: otherwise it is
$(z+1)(z^2+z+1)$. For odd $v$, it splits precisely when $t$ is even:
otherwise it is $z(z^2+z+1)$. The quadratic $z^2+z+1$ is irreducible.
Thus $t+v$ must be odd. Since
\[
 a+c+4b=16(t+v)+11,
\]
this forces the asserted congruence. Its failure gives a noninteger
complement root and graph lift.
\end{proof}

This excludes half the classes left by Proposition~\ref{thm:odd-root-sum}.
There are $512$ role-ordered classes modulo thirty-two; the earlier
condition leaves $128$, and the new one leaves $64$. Thus $64$ classes,
or $384$ after all permutations, are additionally excluded.
For example, the ordered triples $(13,18,19)$ and $(9,31,38)$ satisfy
the old sum condition and $64\mid h_C(3)$, but their residue-role weighted
sums are eleven modulo thirty-two. The compatible binary cubics are
$z^3$ and $(z+1)^3$, so hypothetical integer odd roots would all be
$b$ modulo eight, also requiring $512\mid h_C(b)$.
The surviving congruence classes do not assert integral spectra.

\subsection{Value, derivative and the third shifted cubic}

\begin{proof}[Proof of Proposition~\ref{thm:clustered-odd-roots}, arithmetic conditions]
The proof of Proposition~\ref{thm:odd-root-lift} forces the three hypothetical
integer odd complement roots to be $b$ modulo eight. The two remaining
roots are even. Thus $512\mid h_C(b)$ and $64\mid h_C'(b)$, since each
first-derivative term retains at least two factors divisible by eight.
The exact quintic identity
\[
 h_C(b)=a^2b^2c^2\bigl(a+c-b(b+2)\bigr)
\]
has prefactor with $2$-adic valuation exactly two. It follows that
$128\mid a+c-b(b+2)$, including when the value is zero.

Write $(a,b,c)=(4u+1,4w+3,b(b+2)+128t-a)$.
Coefficientwise expansion gives
\[
 h_C'(b)\equiv16(w^2-w+2)\pmod{64}.
\]
For $w\equiv0,1,2,3\pmod4$, the derivative residues are respectively
$32,32,0,0$. Its required divisibility forces $w\equiv2$ or $3\pmod4$,
which is the stated restriction on $b$.

For the two remaining residues, put $K(z)=h_C(8z+b)/512$.
Coefficientwise identities give $K\in\mathbb Z[u,v,t,z]$ and
\[
 K(z)\equiv
 \begin{cases}
 z^3+z^2+(1+v+u^2)z+t,&b=16v+11,\\
 z^3+z^2+(1+v)z+t,&b=16v+15
 \end{cases}\pmod2.
\]
Writing the odd roots as $b+8r_i$ and the even roots as $\eta_j$ gives
\[
 K(z)=\prod_{j=1}^2(8z+b-\eta_j)\prod_{i=1}^3(z-r_i).
\]
The even-root factors reduce to one, so the displayed binary cubic must
split completely. For $\lambda,t\in\mathbb F_2$, the cubic
$z^3+z^2+\lambda z+t$ splits precisely when $t=\lambda$: the two
possibilities are $z^2(z+1)$ and $(z+1)^3$. Otherwise it is
$z(z^2+z+1)$ or the irreducible cubic $z^3+z^2+1$.
Using $u^2\equiv u\pmod2$ gives the two stated parity conditions.
\end{proof}

\section{Endpoint polynomial proofs, certificates and covering scope}
\label{sec:endpoint-details}

The two possible low integer roots have different arithmetic constraints.
The endpoint two is confined to a linear region in the minimum exponent;
this gives a new nonintegral tail for even exponent triples.
The shifted odd-root distributions extend the same tail to two mixed-parity
patterns below.

\begin{proposition}
\label{thm:endpoint-two-bound}
Let $4\leq a\leq b\leq c$, and put
\[
R(a)=\frac{7a^2-2a+8}{a+2}=7a-16+\frac{40}{a+2}.
\]
If $h_C(2)=0$, then $c<R(a)$. In fact, $c\geq R(a)$ implies
$h_C(2)>0$.
\end{proposition}

\begin{proof}
Fix $a,c$ and expand the symmetric endpoint polynomial in $b$:
\[
h_C(2)=A(a,c)b^3+B(a,c)b^2+D(a,c)b+E(a,c),
\]
where
\begin{align*}
A(a,c)&=ac(a+c-1)+2a(a-1)+2c(c-1)-4,\\
B(a,c)&=c^2\bigl((a+2)c-(7a^2-2a+8)\bigr)\\
&\quad +(a^3+2a^2-6a-4)c+2(a-2)(a^2-2a-6),\\
D(a,c)&=(a-2)(c-2)\bigl(a^2c+a^2+ac^2+6ac+4a+c^2+4c-12\bigr),\\
E(a,c)&=2(a-2)(c-2)(a+c-2)(ac+a+c-2).
\end{align*}
For $a,c\geq4$, the coefficients $A,D,E$ are strictly positive.
Both $a^3+2a^2-6a-4$ and $a^2-2a-6$ are positive at $a=4$
and increasing on $[4,\infty)$. Thus $c\geq R(a)$ also gives $B>0$,
which proves the strict endpoint sign. Conversely, an endpoint zero
requires $B<0$, so $(a+2)c<7a^2-2a+8$.
\end{proof}

\begin{proposition}
\label{cor:even-linear-tail}
If $4\leq a\leq b\leq c$ are all even and $c\geq R(a)$,
then $G_{p^a q^b r^c}$ is not Laplacian integral.
For minimum exponent at least eight, $c\geq7a-12$ suffices.
\end{proposition}

\begin{proof}
Every entry of $C$ is even, so every integer eigenvalue of $C$ is even
by the rational-root theorem applied to $C/2$. In particular, one is
not an eigenvalue. Proposition~\ref{thm:endpoint-two-bound} excludes two,
and Theorem~\ref{thm:uniform-low-root} provides a noninteger root in
$(0,3)$. Equivalently, $h_C(0)<0<h_C(2)$ directly yields a root in
$(0,2)$, which cannot equal one. Its graph lift is noninteger.
For $a\geq8$, $40/(a+2)\leq4$ gives $R(a)\leq7a-12$.
\end{proof}

The triples $(8,10,44)$, $(10,14,58)$ and $(40,42,266)$ meet this
criterion with gcd two. Their largest exponents are below the earlier
quadratic cutoff and their spans are outside the earlier balanced region.
The already settled repeated triple $(10,10,12)$ has $h_C(2)=0$ and
obeys the strict bound; an integer endpoint does not establish integrality.

\begin{proposition}
\label{cor:mixed-linear-tail}
Let $4\leq m\leq d\leq M$ be the ordered exponents of a triple that is
a permutation of $(1,3,2)$ or $(3,3,0)$ modulo four.
If $M\geq R(m)$, the ideal intersection graph is nonintegral.
For $m\geq8$, the simpler bound $M\geq7m-12$ suffices.
\end{proposition}

\begin{proof}
Under the integer-spectrum assumption,
Lemma~\ref{lem:mixed-endpoint-one} excludes one as a quotient root.
Theorem~\ref{thm:uniform-low-root} then forces two to be a root.
But Proposition~\ref{thm:endpoint-two-bound} gives $h_C(2)>0$ for
$M\geq R(m)$, a contradiction. The simpler bound follows from
$R(m)\leq7m-12$ for $m\geq8$.
\end{proof}

Thus the linear upper bound also restricts these two mixed-parity patterns.
For example, the residue-role triples $(25,27,758)$ and $(29,27,882)$
satisfy the arithmetic conditions of Proposition~\ref{thm:clustered-odd-roots},
but their largest exponents exceed this linear bound. The root-distribution
argument therefore settles these compatible controls without a further
congruence lift. The triple $(19,27,140)$ illustrates the other pattern.
The variables $m,d,M$ identify size order, independently of residue roles.

The size bound confines one surface geometrically. The following
constant-term conditions give finite candidate sets for either surface
when two exponents are fixed.

\begin{proposition}
\label{prop:endpoint-divisors}
For fixed $a,b\geq4$, any positive integer $c$ on an endpoint-zero
surface satisfies
\begin{align*}
h_C(1)=0&\ \Longrightarrow\ c\mid(a-1)(b-1)(a+b-1)(ab+a+b-1),\\
h_C(2)=0&\ \Longrightarrow\ c\mid2(a-2)(b-2)(a+b-2)(ab+a+b-2).
\end{align*}
\end{proposition}

\begin{proof}
Both endpoints are integer cubics in $c$. Their constant terms are
the displayed nonzero products. Reduce each endpoint equation modulo $c$.
\end{proof}

For a fixed pair this gives finite divisor candidates, which must still
satisfy the endpoint equation. Reza's requested diagnostic for middle
exponents through fifteen now has a complete independent certificate.

\begin{proposition}
\label{cor:middle-fifteen}
Every positive ordered three-prime exponent triple $a\leq b\leq c$
with $b\leq15$ gives a nonintegral ideal intersection graph.
This conclusion uses the complete finite endpoint certificate below.
\end{proposition}

\begin{proof}
Repeated exponents are covered by Theorem~\ref{thm:repeated}, and
minimum exponents one, two and three by Theorems~\ref{thm:unit},
\ref{thm:minimum-two} and~\ref{thm:minimum-three}, respectively.
It remains to consider $4\leq a<b\leq15$ and $c>b$.

For each of these $66$ pairs, both endpoint polynomials are integer
cubics in $c$ with the positive constant terms in
Proposition~\ref{prop:endpoint-divisors}. The
\href{https://github.com/the-omega-institute/ideal-intersection-laplacian/blob/10ff01f0625c7264b8bc5bb33b1a9d7dcb88a93c/results/endpoint-surfaces-verification.json}{complete endpoint certificate}
records all $132$
coefficient lists and exhausts their $8\,527$ positive divisors using
integer trial division. Independent integer Horner evaluation at every
one of the $7\,585$ divisors above $b$ gives no zero. Thus neither
endpoint vanishes for any integer $c>b$, with no imposed upper bound
on $c$. Theorem~\ref{thm:uniform-low-root} supplies a positive root
in $(0,3)$; it can equal neither one nor two, proving nonintegrality.
\end{proof}

The reflected quotient in this diagnostic uses
$N=|V(H)|=a+b+c+ab+ac+bc$, not the full graph vertex count.
The inclusive requested bound $b\leq15$ and exact rational linear-factor
extraction are cross-checked against the independent divisor certificate.
The empty lists settle this pair range; they do not prove either
endpoint surface globally empty.

\begin{proposition}
\label{thm:modulo-three-endpoints}
Every positive three-prime exponent triple congruent to $(2,2,2)$,
or to a permutation of $(1,2,2)$, modulo three gives a nonintegral
ideal intersection graph.
\end{proposition}

\begin{proof}
For these two residue patterns, reduction of the general complement
quintic gives, respectively,
\[
h_C(x)\equiv x^5,\qquad
h_C(x)\equiv(x^2+1)(x^3+x^2-x+1)\pmod3.
\]
The corresponding pairs $(h_C(1),h_C(2))$ are $(1,2)$ and $(1,1)$
modulo three, so neither integer endpoint is a root. Minimum exponents
at most three are already covered. Otherwise
Theorem~\ref{thm:uniform-low-root} supplies a positive root in $(0,3)$,
which cannot be an integer. Its graph lift is noninteger as well.
For the second pattern, the irreducible factor $x^2+1$ also directly
rules out a completely integral quotient spectrum.
\end{proof}

For example, $(8,17,22)$ and $(8,17,23)$ satisfy the new conditions
with gcd one and unequal $2$-adic valuations. Complete endpoint tables
modulo two and three, checked against independent integer determinants,
are given in the \href{https://github.com/the-omega-institute/ideal-intersection-laplacian/blob/2df3cffd062ce0b28fc2040a6203fd4e7282a725/results/endpoint-residues.json}{residue certificate}.
Every pair $(a,b)$ modulo either prime has at least one $c$ root on
each surface; retaining the full root sets, rather than only root-free
pairs, is necessary to detect these exclusions.

\begin{remark}
Endpoint tests at a fixed finite set of moduli cannot by themselves
prove either surface empty on all fully distinct positive triples.
Let $M$ be their least common multiple. Since $(9,9,136)$ has
$h_C(1)=0$, the strictly ordered triples
$(9+kM,9+(k+1)M,136+(k+2)M)$ pass every selected endpoint-one test.
The known endpoint-two zero $(10,10,12)$ similarly gives survivors
$(10+kM,10+(k+1)M,12+(k+2)M)$. For large $k$ their middle exponents
exceed fifteen. These are congruence survivors, not asserted integer
zeros. These fixed-span tails are excluded by
Corollary~\ref{cor:balanced-span}. To retain the geometric restrictions,
take $L$ divisible by $M$ and eight and any $d\geq32$ divisible by $L$.
The triples $(9+d,9+2d,136+4d)$ and $(10+d,10+2d,12+4d)$ pass the
respective endpoint tests, have gcd one and two, and satisfy both upper
bounds on the maximum while failing the sufficient inertia inequality.
The substitution identities are given in the accompanying
\href{https://github.com/the-omega-institute/ideal-intersection-laplacian/blob/nikandish-math-patch-1/notes/endpoint-congruence-scope.md}{scope note}.
This concerns endpoint-value tests alone; another-root obstruction or
further arithmetic information can still exclude these candidates.
\end{remark}

Theorem~\ref{thm:endpoint-two-completion} completes the endpoint-two
alternative, including the all-even triples formerly left inside this
linear region. The remaining mixed-parity surface $h_C(1)=0$ still
needs another obstruction.

\end{document}
